% Validation: 56 source pages; 27 numbered equations; 22 figures; 31 numbered footnotes; 63 references.
% Mathematical displays checked against rendered source; prose coverage checked across all 56 pages.
% Reflowed compilation inspected; figures depend on the adjacent source PDF.
% Source PDF SHA-256: 4a238406459a58d2a129cfee5ade258a1946cb848f361fa33ba505a86a425738
% Agent-readable reference transcription of the supplied published PDF.
% This is a PDF reconstruction, not the authors' original TeX source.
% Compile from this directory with: pdflatex Atkeson_2025_End_of_Privilege.tex
% Keep the original PDF beside this file: figures are vector crops from that PDF.
% Original numbered equations: eq:1 through eq:27.
% Original page anchors: source-pdf-1 through source-pdf-56 (journal 2151--2206).
% All editorial remarks are explicitly marked Conversion note.
\documentclass[11pt]{article}
\usepackage[T1]{fontenc}
\usepackage[utf8]{inputenc}
\usepackage{lmodern,amsmath,amssymb,graphicx,booktabs,array,geometry,xcolor}
\usepackage[unicode,hidelinks]{hyperref}
\usepackage{bookmark}
\geometry{margin=24mm}
\allowdisplaybreaks
\setlength{\emergencystretch}{3em}
\setlength{\parindent}{0pt}
\setlength{\parskip}{6pt}
\newcommand{\sourcepage}[2]{\clearpage\hypertarget{source-pdf-#1}{}\pdfbookmark[2]{Original PDF #1 / journal #2}{page-#1}\noindent{\small\color{gray}Original PDF page #1; journal page #2}\par\medskip}
\newcommand{\sourcepdf}{\detokenize{Atkeson et al. - 2025 - The End of Privilege A Reexamination of the Net Foreign Asset Position of the United States.pdf}}
\newcommand{\sourcedisplay}[1]{\par\begingroup\small #1\par\endgroup}
\begin{document}
\begin{center}
{\LARGE The End of Privilege: A Reexamination of the Net Foreign Asset Position of the United States\par}
\medskip
Andrew Atkeson, Jonathan Heathcote, and Fabrizio Perri\\
\textit{American Economic Review} 115(7), July 2025, pp. 2151--2206\\
\url{https://doi.org/10.1257/aer.20230732}
\end{center}
\paragraph{Conversion note.} This reference transcribes the supplied 56-page published article. It retains the original equation numbers, figure and table numbers, footnote numbers, and journal page references. Original page boundaries are explicit; a paragraph may therefore continue on the next source page. Figures are vector crops of the adjacent original PDF. This file has been reflowed for reading, so its compiled pagination differs from the journal. The separate Supplemental Appendix mentioned by the authors is not contained in the supplied PDF.
\paragraph{Source fidelity.} Mathematical expressions are transcribed as printed, including apparent typographical inconsistencies; conversion comments near the affected equations identify conspicuous cases. They should not be interpreted as corrections to the authors' model. Source annotations and highlights are not part of the article text.
\tableofcontents


% ===== ORIGINAL PDF PAGE 1; JOURNAL PAGE 2151 =====
\sourcepage{1}{2151}

% Source PDF 1, text block 4.
\begin{abstract}
The US net foreign asset position declined sharply after 2007, reaching negative 60 percent of GDP by the third quarter of 2023. This deterioration primarily reflects a US-specific rise in corporate asset values that inflated the value of US equity liabilities to the rest of the world. To interpret these trends, we develop an international macrofinance model of flows, stocks, asset valuations, the current account, and the net foreign asset position. We find that the welfare impact of rising asset values for a representative US household has been quite negative given extensive foreign ownership of US corporate equity. (JEL F32, F21, F40, G15)
\end{abstract}


% Source PDF 1, text block 5.
Figure 1 plots the net foreign asset position and current account of the United States from 1990 to the third quarter of 2023. The net foreign asset position (henceforth, NFA) is measured as the market value of the assets US residents hold abroad minus the market value of US assets held by residents of the rest of the world. The figure shows that over the 1990–2007 period, the United States maintained a relatively small negative net position, despite running sustained and substantial current account deficits. As discussed by Gourinchas and Rey (2007) and Gourinchas, Rey, and Govillot (2017), up until 2007, US residents enjoyed the privilege of being able to borrow from the rest of the world without increasing US net debt thanks to ex post favorable market revaluations of cross-border assets and liabilities.


% Source PDF 1, text block 6.
In sharp contrast to this prior experience, from 2007 into 2021 the US NFA position declined precipitously—by 60 percentage points of US GDP—before bouncing back somewhat in 2022. And this occurred despite the fact that US current account deficits narrowed relative to the early 2000s.


% Source PDF 1, text block 7.
We document that this unprecedented decline in the US NFA position was driven by a boom in the market valuation of the nonfinancial assets in US corporations. Because foreigners’ gross holdings of equity in US corporations had grown to be


% Source PDF 1, text block 8.
* Atkeson: Department of Economics, UCLA (email: andy.atkeson@gmail.com); Heathcote: Research Department, Federal Reserve Bank of Minneapolis (email: heathcote@minneapolisfed.org); Perri: Research Department, Federal Reserve Bank of Minneapolis (email: fabri.perri@gmail.com). Arnaud Costinot was the coeditor for this article. This is an extensively revised version of our 2022 paper with the same title. We thank our editor and three referees for detailed and constructive suggestions. We are grateful to Robert Barro, Gianluca Benigno, Carol Bertaut, John Cochrane, Stephanie Curcuru, Mick Devereux, Pierre-Olivier Gourinchas, Dan Greenwald, Şebnen Kalemli-Özcan, Hanno Lustig, Matteo Maggiori, Thomas Philippon, Vincenzo Quadrini, Hélène Rey, and Jesse Schreger for helpful comments. Thanks to Jiaxi Tan for excellent research assistance. The views expressed herein are those of the authors and not necessarily those of the Federal Reserve Bank of Minneapolis or the Federal Reserve System.


% Source PDF 1, text block 9.
\textdagger{} Go to https://doi.org/10.1257/aer.20230732 to visit the article page for additional materials and author disclosure statement(s).



% ===== ORIGINAL PDF PAGE 2; JOURNAL PAGE 2152 =====
\sourcepage{2}{2152}

% Figure 1: original PDF page 2, plot crop x=60..444, y=62..228 points.
\begin{center}
\begingroup\renewcommand{\thefigure}{1}\refstepcounter{figure}\label{fig:1}\endgroup
\includegraphics[page=2,trim=60 492 60 62,clip,width=0.95\linewidth]{\sourcepdf}
\par\small Figure 1. The US Net Foreign Asset Position and Current Account, 1990–2023
\end{center}


% Source PDF 2, text block 1.
very large, this boom mechanically increased the market value of US liabilities to the rest of the world (henceforth, ROW). There was no similar boom in the valuation of corporations in the ROW over this time period, so US residents did not enjoy a similar revaluation of their gross foreign equity assets. As a result, the net impact of asset revaluations accounted for a large portion of the deterioration in the US NFA position between the Great Recession and 2023. In fact, as we show below, the negative impact of these revaluations of gross cross-border equity positions was so large that the US NFA position was more negative by the third quarter of 2023 than it would have been had no asset revaluations occurred between 1990 and 2023. In this sense, the ex post “privilege” that US residents previously enjoyed was erased.\hyperlink{fn:1}{\textsuperscript{1}}


% Source PDF 2, text block 2.
Motivated by these observations, we ask two questions. First, what factors underlie this deterioration of the US NFA position and the boom in the market valuation of US corporations? Second, what do these developments mean for the welfare of US residents?


% Source PDF 2, text block 3.
To answer these two questions, we develop a unified international macrofinance model of flows, stocks, and valuations of the US corporate sector and of the US current account and NFA position. The model builds on Farhi and Gourio (2018) and Crouzet and Eberly (2023) but extends those frameworks to an international setting to include international positions and flows.\hyperlink{fn:2}{\textsuperscript{2}} The approach we take to integrating the current account follows the model of intertemporal trade under incomplete markets of Obstfeld and Rogoff (1995); Engel and Rogers (2006); Corsetti


% Source PDF 2, text block 4.
{\footnotesize \hypertarget{fn:1}{}\textbf{Footnote 1.} Gourinchas and Rey (2014); Gourinchas, Rey, and Govillot (2017); Chen et al. (2022); Choi, Kirpalani, and Perez (forthcoming); Gourinchas (2023); and many others discuss the role of ex ante return differentials on US foreign assets and liabilities in shaping the US external position. See Curcuru, Thomas, and Warnock (2013) and Bertaut et al. (2024) for critical reviews of the evidence for an ex ante difference in expected returns on US foreign assets and liabilities. In our analysis, we do not assume any ex ante return differential on US assets and liabilities. We focus our analysis on the impact of differences in ex post realized returns. Other authors have also highlighted the large boom in the value of US assets and its impact on the US NFA position; see, for example, Jiang, Richmond, and Zhang (2024); Milesi-Ferretti (2021, 2023).\par}


% Source PDF 2, text block 5.
{\footnotesize \hypertarget{fn:2}{}\textbf{Footnote 2.} See also Caballero, Farhi, and Gourinchas (2017); Gutiérrez and Philippon (2017); Eggertsson, Robbins, and Wold (2021); Greenwald, Lettau, and Ludvigson (2025); Corhay, Kung, and Schmid (2025); and others for macrofinance models of the boom in the valuation of US corporations.\par}



% ===== ORIGINAL PDF PAGE 3; JOURNAL PAGE 2153 =====
\sourcepage{3}{2153}

% Source PDF 3, text block 1.
and Konstantinou (2012); and many others. Model households in two regions (the United States and ROW) trade domestic and foreign equity and risk-free bonds. Firms in both countries enjoy pricing power that translates to rents payable to their shareholders that Karabarbounis and Neiman (2019) refer to as factorless income. The size of this factorless income can vary across countries and over time, generating fluctuations in equity valuations relative to value added. Additional sources of time variation in asset values include fluctuations in the equilibrium discount rate applied to future cash flows, fluctuations in expected future growth rates, fluctuations in the replacement cost of capital, and fluctuations in corporate tax rates.


% Source PDF 3, text block 2.
The model is fully tractable. We exploit its tractability to measure the factors driving observed flows, stocks, and valuations of the US corporate sector, together with the evolution of the US current account and NFA position, in quarterly data over the period 1990–2023. We saturate the model with interpretable time-varying parameters and compute sequences for parameter values such that the model matches the data period by period. We then simulate counterfactual model scenarios relative to our baseline to study how the factors driving equity values impacted the welfare of US households. We have two main findings.


% Source PDF 3, text block 3.
First, we find that much of the increase in the market valuation of the nonfinancial assets in US corporations since the Great Recession has been due to a dramatic increase in the free cash flow from operations available to pay to owners of firms that is unprecedented in post-WWII data. This measure of corporate income is defined directly from the national accounts as the amount left over from corporate sector value added after deducting payments to labor, taxes (both indirect business taxes and taxes on corporate profits), and investment expenditures on new nonfinancial assets.\hyperlink{fn:3}{\textsuperscript{3}} We find that changes in the valuation multiple applied to free cash flow have played a much smaller role in driving the increased valuation of US corporations. In our accounting, some of this increase in corporate free cash flow is due to changes in taxes and the share of labor in costs, but the lion’s share is due to an increase in the wedge between revenue and total cost, resulting in a large increase in the share of factorless income in US corporate gross value added. This finding is consistent with earlier work by Farhi and Gourio (2018); Barkai (2020); Eggertsson, Robbins, and Wold (2021); and, most relatedly, Greenwald, Lettau, and Ludvigson (2025). In what follows, we refer to the wedge between revenue and total cost as the output wedge.\hyperlink{fn:4}{\textsuperscript{4}}


% Source PDF 3, text block 4.
Second, when we use our model to simulate counterfactuals, we find that the welfare implications of these developments driving the increase in valuation of US corporations are dramatically impacted by the observed large increase in gross cross-border equity positions. Specifically, we find that had US residents been the sole owners of US corporations, the observed rise in the output wedge would have had only a small impact on the welfare of a representative US household. This


% Source PDF 3, text block 5.
{\footnotesize \hypertarget{fn:3}{}\textbf{Footnote 3.} In contrast, related measures of corporate income, such as estimates of pure profits (e.g., Barkai 2020) or of corporate earnings (e.g., Greenwald, Lettau, and Ludvigson 2025) require imputations of compensation to physical capital or estimates of depreciation.\par}


% Source PDF 3, text block 6.
{\footnotesize \hypertarget{fn:4}{}\textbf{Footnote 4.} We use the terminology output wedge rather than markup to emphasize that this wedge measures pure profits, that is, the gap between price and the average cost of production factors. It plays the same role as the output distortion in Hsieh and Klenow (2009) and the “markups” in Farhi and Gourio (2018); Baqaee and Farhi (2020); Barkai (2020); and Crouzet and Eberly (2018). It is distinct from the measure of markups of price over marginal variable cost, which is the focus of, for example, De Loecker, Eeckhout, and Unger (2020).\par}



% ===== ORIGINAL PDF PAGE 4; JOURNAL PAGE 2154 =====
\sourcepage{4}{2154}

% Source PDF 4, text block 1.
welfare impact would have been small because lower wage income would have been largely offset by higher free cash flow to US households as owners of US corporations.\hyperlink{fn:5}{\textsuperscript{5}} In contrast, given the large cross-border equity positions observed in the data, we find that the observed rise in the output wedge has a large negative impact on the consumption of US residents. The reason is simple: Much of the increase in free cash flow of US firms is paid to foreign owners.


% Source PDF 4, text block 2.
Could the transfer of resources from US residents to foreign equity owners that follows a rise in the US output wedge be interpreted as part of an ex ante efficient international risk-sharing arrangement? We explore this question in an extended version of the model and show that the extent of international portfolio diversification that delivers maximal risk sharing is highly sensitive to the nature of country-specific risk. If that risk takes the form of shocks to the US output wedge, then the exposure of US residents to domestic equity appears close to optimal.


% Source PDF 4, text block 3.
We make three principal contributions to the literature.

First, we build a model to provide an integrated accounting of flows, stocks, and valuations of the US corporate sector and of the US current account and NFA position. It has long been recognized that the current account and net foreign asset position of a country are impacted not only by changes in capital accumulation but also through changes in asset valuations, both directly through revaluations of existing cross-border asset holdings and indirectly through wealth effects impacting the ratio of consumption to income.\hyperlink{fn:6}{\textsuperscript{6}} While all of these effects are present qualitatively in standard international business cycle models, these standard models typically do not account quantitatively for the large changes in valuations of firms at home and abroad observed in the data. Here, we address this shortcoming of standard international business cycle models by extending the recent macrofinance literature that has been developed to account for large observed changes in the valuation of US corporations. We use this model to better understand the links between changes in asset valuations on the one hand and current account and NFA dynamics on the other.


% Source PDF 4, text block 4.
Second, we bring additional data to bear on the question of whether the observed increase in the market valuation of the US corporate sector is driven by an increase in cash flows to owners of firms or by a change in the valuation multiple of those cash flows. By including data on the current account in our measurement, we are able, through the model, to separately identify the impact of these factors. Our model-based measurement comes down in favor of a stable ratio of expected free cash flow to the market value of nonfinancial assets in US corporations over the past decade, as the model requires a relatively stable valuation multiple to account for the relative stability of the US current account balance.


% Source PDF 4, text block 5.
Third, and perhaps most important, as we discuss above, we find that conclusions regarding the welfare costs to US residents of a large increase in the share of corporate value added attributed to factorless income are highly sensitive to the extent of cross-border diversification of equity positions.


% Source PDF 4, text block 6.
{\footnotesize \hypertarget{fn:5}{}\textbf{Footnote 5.} See, for example, Corollary 1 in Baqaee and Farhi (2020) for a theoretical derivation of this result for a small change in the output wedge in a closed economy. This quantitative finding in our model is obtained given a large increase in the output wedge.\par}


% Source PDF 4, text block 7.
{\footnotesize \hypertarget{fn:6}{}\textbf{Footnote 6.} See, for example, Lane and Milesi-Ferretti (2001, 2007) and subsequent updates of these data by these authors, and the discussion of the literature on the current account and NFA position in Gourinchas and Rey (2014).\par}



% ===== ORIGINAL PDF PAGE 5; JOURNAL PAGE 2155 =====
\sourcepage{5}{2155}

% Source PDF 5, text block 1.
One important assumption in our baseline analysis is that the valuation multiple relevant for future factorless income is also the one used to value future labor income. The representative US household in our model can borrow against future labor income, and changes in the discount factor or the long-term expected growth rate therefore impact the current account via their impact on US household human wealth. Lustig and Van Nieuwerburgh (2008) show that this assumption in conjunction with human wealth returns that are negatively correlated with returns to financial wealth can explain the joint behavior of consumption and returns to financial wealth in the United States. However, other studies (see, for example, Greenwald, Lettau, and Ludvigson 2025) explain the joint behavior of consumption and asset prices in a setup where a large fraction of households are not able to borrow against future labor income. For this reason, we assess the sensitivity of our findings to this assumption, by conducting an array of alternative measurement exercises in which we do not use data on the current account to identify parameter values.


% Source PDF 5, text block 2.
Instead, we use a measurement procedure similar to that in Farhi and Gourio (2018) and Crouzet and Eberly (2023) to consider a wide range of alternative scenarios for expected future growth rates or for the gap between the discount factor and expected future growth.\hyperlink{fn:7}{\textsuperscript{7}} We find that the conclusion that much of the increase in the market valuation of US corporations over the past decade is due to an increase in the output wedge is robust to this wide range of alternative measurements.


% Source PDF 5, text block 3.
The remainder of the paper is organized as follows. In Section I and Appendix A, we present the data we use on the evolution of the US NFA position and current account since 1990, and the flows, stocks, and valuation of the US corporate sector. In Section II, we describe the international macrofinance model we use to interpret this data. In Section III and Section 2 of the Supplemental Appendix, we describe how we use the model to measure the factors driving the flows, stocks, and valuation of the US corporate sector as well as the US current account and NFA position. We present our baseline findings in Section IV. Section V contains our counterfactual exercises to evaluate the welfare impact of changing valuations. Sensitivity analysis is in Section VI.


% Source PDF 5, text block 4.
\section*{I. The Evolution of the US Current Account, NFA Position, and the US Corporate Sector: 1990–2023}
\addcontentsline{toc}{section}{I. The Evolution of the US Current Account, NFA Position, and the US Corporate Sector: 1990–2023}


% Source PDF 5, text block 6.
In this section, we review the measurement concepts and data we analyze with our model. We begin with a discussion of the evolution of the US NFA position and then turn to the data on flows, stocks, and valuation of the US corporate sector. Details on the data series used are provided in Appendix A.


% Source PDF 5, text block 7.
The NFA Position and Its Components.—The starting point of our analysis is accounting identity (1), showing that the change in the NFA position between the end of periods $t-1$ and $t$ is the sum of three components. The first, $(CA_{t})$, is the balance of the current account during period $t$; this term captures net US lending abroad measured as the sum of net exports and net income receipts. The second


% Source PDF 5, text block 8.
{\footnotesize \hypertarget{fn:7}{}\textbf{Footnote 7.} We describe how our measurement procedure relates to that used in prior macrofinance papers in greater detail in Section 3 of the Supplemental Appendix.\par}



% ===== ORIGINAL PDF PAGE 6; JOURNAL PAGE 2156 =====
\sourcepage{6}{2156}

% Source PDF 6, text block 1.
term, $(VA_{t})$, captures the net change in the valuations of the existing assets that compose the gross external positions. The third term, $(RES_{t})$, is a residual, which reconciles the changes in the NFA position resulting from measured financial transactions and asset positions with the ones resulting from current account transactions.\hyperlink{fn:8}{\textsuperscript{8}} Thus,


% Source PDF 6, text block 2.
\setcounter{equation}{0}
\begin{equation}\label{eq:1}
NFA_t-NFA_{t-1}=\underbrace{CA_t}_{\text{net lending abroad}}+\underbrace{VA_t}_{\text{valuation effects}}+\underbrace{RES_t}_{\text{residual term}}.
\end{equation}


% Source PDF 6, text block 15.
Summing (1) from period 1 to period $t$ yields


% Source PDF 6, text block 16.
\setcounter{equation}{1}
\begin{equation}\label{eq:2}
NFA_t=NFA_0+\underbrace{\sum_{j=1}^{t}CA_j}_{\text{cumulated CA}}+\underbrace{\sum_{j=1}^{t}VA_j}_{\text{cumulated valuations}}+\underbrace{\sum_{j=1}^{t}RES_j}_{\text{cumulated residuals}},
\end{equation}


% Source PDF 6, text block 38.
showing that the NFA position in any period can be expressed as the cumulated sums of the three terms described above.


% Source PDF 6, text block 39.
Figure 2 shows the evolution of the three components in equation (2), each divided by US corporate gross value added (GVA) in quarter $t$, from 1990:I until 2023:III. The figure shows three different phases in the evolution of the US NFA position. During the first phase (1990–2002), the NFA position closely tracked cumulative current account dynamics. During the second phase (2002–2007), the cumulative current account continued to deteriorate, but the NFA position improved, owing to a combination of positive valuation effects and positive statistical discrepancies. This period was the focus of Gourinchas and Rey (2007, 2014), who noticed that valuation effects, which increased the value of foreign assets held by US residents relative to the value of US assets held by foreigners, acted as a stabilizing counterweight to growing current account deficits. In the third and final phase (2007–2022), the US NFA position declined substantially, despite a fairly stable (relative to corporate GVA) cumulated current account deficit. Note that by 2020, the US NFA position was more negative than cumulated current accounts over the entire 1990 to 2020 period. As is evident in the figure, a large portion of the decline of the US NFA position in this third phase was driven by negative valuation effects, meaning that during this period, US residents experienced consistently lower capital gains on their foreign asset holdings than those enjoyed by ROW residents on their US assets.\hyperlink{fn:9}{\textsuperscript{9}}


% Source PDF 6, text block 40.
Decomposing Valuation Effects.—Since cumulated valuation effects are an important determinant of the evolution of the US NFA position, we now proceed to analyze in more detail the sources and the impacts of these valuation changes. As a matter of accounting, valuation effects are given by


% Source PDF 6, text block 41.
\setcounter{equation}{2}
\begin{equation}\label{eq:3}
VA_t=FA_{t-1}\times g_t^{P^*}-FL_{t-1}\times g_t^P,
\end{equation}


% Source PDF 6, text block 42.
{\footnotesize \hypertarget{fn:8}{}\textbf{Footnote 8.} See Curcuru, Dvorak, and Warnock (2008, Section 3) for a discussion of these discrepancies arising from differences in the measurement of international financial flows and positions.\par}


% Source PDF 6, text block 43.
{\footnotesize \hypertarget{fn:9}{}\textbf{Footnote 9.} Figure A1 in Appendix A shows that these patterns are also evident in an alternative version of this decomposition using the cumulated current account from financial transactions.\par}



% ===== ORIGINAL PDF PAGE 7; JOURNAL PAGE 2157 =====
\sourcepage{7}{2157}

% Figure 2: original PDF page 7, plot crop x=60..444, y=62..290 points.
\begin{center}
\begingroup\renewcommand{\thefigure}{2}\refstepcounter{figure}\label{fig:2}\endgroup
\includegraphics[page=7,trim=60 430 60 62,clip,width=0.95\linewidth]{\sourcepdf}
\par\small Figure 2. Decomposition of Changes in the US Net Foreign Asset Position as a Share of US Corporate Value Added
\end{center}


% Source PDF 7, text block 1.
where $FA_{t-1}$ and $FL_{t-1}$ are gross US net foreign asset and liability positions at the end of $t-1$ and $g_{t}^{P^{*}}$ and $g_{t}^{P}$ are the net growth rates in the dollar values of those positions between the end of $t-1$ and the end of period $t$. It is immediate from this expression that there are two necessary conditions for valuation effects to matter quantitatively: (i) gross positions must be large, and (ii) the values of foreign assets and foreign liabilities cannot co-move too closely. We now document that both these conditions have been satisfied in the past decade.


% Source PDF 7, text block 2.
It is useful to divide US foreign positions into two broad categories: equity and non-equity investments. Equity investment includes portfolio investment in corporate equities and the equity component of foreign direct investment (FDI). At the beginning of our sample, when international equity markets were still relatively underdeveloped, FDI was the main component of both inward and outward equity investment, accounting for 80 percent of both positions. Toward the end of our sample, with large and active international equity markets, portfolio and direct equity investment have roughly equal shares. Non-equity assets include debt securities, loans, and currency and deposits. Figure 3 plots the evolutions of these categories of US foreign assets and liabilities as fractions of US corporate GVA.


% Source PDF 7, text block 3.
The first key message from Figure 3 is that by 2007, all the gross positions are large, and thus, changes in the prices of the assets composing these positions can potentially generate significant valuation effects. The second key message is that US equity liabilities have grown dramatically since the early 1990s, and US equity liabilities now exceed US equity foreign assets. Thus, changes in the price of US equity that are not matched by identical changes in the price of ROW equity now have much larger effects on the US NFA position than would have been the case in the past.



% ===== ORIGINAL PDF PAGE 8; JOURNAL PAGE 2158 =====
\sourcepage{8}{2158}

% Figure 3: original PDF page 8, plot crop x=60..444, y=62..227 points.
\begin{center}
\begingroup\renewcommand{\thefigure}{3}\refstepcounter{figure}\label{fig:3}\endgroup
\includegraphics[page=8,trim=60 493 60 62,clip,width=0.95\linewidth]{\sourcepdf}
\par\small Figure 3. Gross Equity and Non-equity Positions over US Corporate Value Added
\end{center}


% Source PDF 8, text block 1.
We now turn to changes in asset valuations. The left panel of Figure 4 decomposes the cumulated valuation effects plotted in Figure  2 into valuation effects arising from equity and non-equity positions. The figure shows that net valuation changes arise almost exclusively from the equity positions. Although in principle both categories are subject to relative valuation changes (due both to price changes and to exchange rate movements for assets denominated in different currencies), these effects are quantitatively much more important for the equity positions.\hyperlink{fn:10}{\textsuperscript{10}}


% Source PDF 8, text block 2.
Why are equity valuation effects so large? The right panel of Figure 4 plots the price indexes, both in dollar terms, that the Bureau of Economic Analysis uses to revalue US equity assets and liabilities. We have normalized the price indexes so that both are equal to one in the first quarter of 2009. Recall that the equity valuation effects in the left panel of Figure 4 are computed by multiplying the gross equity asset and liability positions plotted in Figure 3 by the growth in these price indexes (equation (3)).


% Source PDF 8, text block 3.
In the 1990s, the price of US equity liabilities rose more rapidly than the price of US foreign equity assets. But cumulated valuation effects were small because gross international equity positions were relatively small in the early part of our sample (Figure 3), so international differentials in equity price dynamics did not translate into large effects on the value of the NFA position.


% Source PDF 8, text block 4.
By the mid-2000s, gross cross-border equity positions were larger, and because equity markets in the ROW outperformed the United States during this period, the US NFA position improved.


% Source PDF 8, text block 5.
In the post–Great Recession period, gross equity positions were larger still. A dramatic rise in US equity prices over this period applied to very large gross equity liabilities led to a sharp increase in the value of US equity liabilities. In particular, the dollar price of US equity liabilities peaked in 2021 at over 4.5 times the price at the end of 2008, while the price of US foreign equity assets rose by only a factor of 2. Thus, the


% Source PDF 8, text block 6.
{\footnotesize \hypertarget{fn:10}{}\textbf{Footnote 10.} One reason why valuation effects for non-equity assets are so small is that foreign bonds owned by Americans tend to be dollar denominated, as are bond liabilities (see Maggiori, Neiman, and Schreger 2020). Regarding the equity valuation effects, in Figure 1 in Section 1 of the Supplemental Appendix, we break down the cumulated valuation changes for equity into those coming from FDI equity versus those from portfolio investment in equity; cumulated valuation effects for equity are roughly equally split between the two components.\par}



% ===== ORIGINAL PDF PAGE 9; JOURNAL PAGE 2159 =====
\sourcepage{9}{2159}

% Figure 4: original PDF page 9, plot crop x=60..444, y=62..225 points.
\begin{center}
\begingroup\renewcommand{\thefigure}{4}\refstepcounter{figure}\label{fig:4}\endgroup
\includegraphics[page=9,trim=60 495 60 62,clip,width=0.95\linewidth]{\sourcepdf}
\par\small Figure 4. Valuation Effects and Equity Prices
\end{center}


% Source PDF 9, text block 1.
value of US foreign equity assets rose by much more than the value of equity liabilities, translating into an unprecedented decline in the US NFA position.


% Source PDF 9, text block 2.
To further illustrate how the size of valuation effects depends on the size of gross equity position, Figure 5 compares actual valuation effects to valuation effects under two counterfactuals that hold fixed the size of gross equity positions at a low value (red dotted line) and a high value (red dashed line). In particular, the dotted (dashed) line shows what valuation effects would have been had equity prices evolved exactly as they did—the right panel of Figure 4—but if US gross equity assets and liabilities had been only 20 percent (120 percent) of US corporate GVA in every quarter.\hyperlink{fn:11}{\textsuperscript{11}} The dashed line illustrates that valuation effects would have been much larger in the 1990s and early 2000s had gross cross-country equity positions in those years been as large as they were in the subsequent period. The plot also illustrates that the large negative cumulative valuation effects observed in the post–Great Recession period mechanically reflect differential equity price movements applied to large gross positions.


% Source PDF 9, text block 3.
To summarize, differential equity price growth matters more for the NFA position when gross international equity positions are large. US equity markets outperformed in the post-2010 period, precisely when it mattered most for the NFA position.


% Source PDF 9, text block 4.
Role of Exchange Rate Movements in Equity Valuation Effects.—Do equity valuation changes arise from changes in the local currency price of equity or from movements in exchange rates? In this section we answer this question for the two notable swings in the US net foreign asset position in our sample: the upward swing from 2003 to 2007 and the downward swing from 2010 to 2020 (see Figure 4). In Figure 6 we plot three equity price indexes. The first (labeled MSCI USA) is a price index for equity in the United States; the second and third are price indexes for equity in the rest of the world measured in US dollars (MSCI ROW in USD) and in local currency (MSCI ROW in LOC).\hyperlink{fn:12}{\textsuperscript{12}} Focus first on the left panel, which


% Source PDF 9, text block 5.
{\footnotesize \hypertarget{fn:11}{}\textbf{Footnote 11.} We thank Dan Greenwald for suggesting this plot.\par}

{\footnotesize\hypertarget{fn:12}{}\textbf{Footnote 12.} The indexes (MSCI Inc. 2024) are the Morgan Stanley Capital Index (MSCI) US Price Index and the MSCI All Country World Price Index ex USA, which comprises stock market indexes (including large and midcap stocks) for 22 developed economies and 27 emerging markets, weighted by market capitalization, in US dollars and in local currency. Bertaut et al. (2024) show that the MSCI indexes closely track the valuation changes reported by the\par}



% ===== ORIGINAL PDF PAGE 10; JOURNAL PAGE 2160 =====
\sourcepage{10}{2160}

% Figure 5: original PDF page 10, plot crop x=60..444, y=62..294 points.
\begin{center}
\begingroup\renewcommand{\thefigure}{5}\refstepcounter{figure}\label{fig:5}\endgroup
\includegraphics[page=10,trim=60 426 60 62,clip,width=0.95\linewidth]{\sourcepdf}
\par\small Figure 5. Counterfactual Valuation Effects
\end{center}


% Figure 6: original PDF page 10, plot crop x=60..444, y=329..495 points.
\begin{center}
\begingroup\renewcommand{\thefigure}{6}\refstepcounter{figure}\label{fig:6}\endgroup
\includegraphics[page=10,trim=60 225 60 329,clip,width=0.95\linewidth]{\sourcepdf}
\par\small Figure 6. Two Valuation Episodes
\end{center}


% Source PDF 10, text block 1.
describes the earlier episode. This panel shows that foreign equity performed better than US equity in local currency, but in dollar terms, the foreign equity index substantially outperformed the US index. Depreciation of the US dollar against the basket of currencies that are represented in the rest-of-world equity index accounted for around half of the positive valuation effects experienced by the United States.\hyperlink{fn:13}{\textsuperscript{13}}


% Source PDF 10, text block 2.
{\footnotesize\textbf{Footnote 12 (continued).} Bureau of Economic Analysis in the International Investment Position tables. These indexes are available from the MSCI website; see \url{https://www.msci.com/end-of-day-data-search}.\par}


% Source PDF 10, text block 3.
{\footnotesize \hypertarget{fn:13}{}\textbf{Footnote 13.} See Bureau of Economic Analysis (2014) for more discussion of how valuation changes reflect a mix of changes in the prices of the underlying assets and changes in exchange rates when assets and liabilities are denominated in different currencies. Note that the revaluation of ROW equity in the United States arises solely because of changes in the price of US equity, as these assets are valued in dollars.\par}



% ===== ORIGINAL PDF PAGE 11; JOURNAL PAGE 2161 =====
\sourcepage{11}{2161}

% Source PDF 11, text block 7.
\refstepcounter{table}\label{tab:1}
\begin{center}
\textbf{Table 1---Impact of Exchange Rate Movements on Equity Revaluations}
\end{center}
\begin{center}\small
\begin{tabular}{@{}p{0.72\linewidth}rr@{}}\toprule
 & 2003--2007 & 2010--2020\\\midrule
(1) Revaluation of US equity abroad & 62.6 & 41.3\\
\quad (1a) contribution of exchange rate movements & 15.8 & $-10.1$\\
\quad (1b) contribution of changes in local currency prices & 46.8 & 51.4\\
(2) Revaluation of foreign equity in US & 22.3 & 103.6\\
(3) Revaluation of net US foreign equity = (1)-(2) & 40.4 & $-62.3$\\\bottomrule
\end{tabular}\end{center}



% Source PDF 11, text block 10.
Note: Changes in value expressed as percentages of US corporate gross value added.


% Source PDF 11, text block 11.
Source: Author calculations based on Bureau of Economic Analysis International Investment Position Table 1.3.


% Source PDF 11, text block 1.
Moving now to the right panel, we can see that the later valuation episode is different. During this period, the United States and rest-of-world equity indexes diverge dramatically. Comparing the foreign indexes in local currency and in dollars indicates some appreciation of the US dollar, but this appreciation accounts for only a small portion of the differential in dollar returns. Rather, the dominant factor is that the US equity price index more than tripled over the period, while the rest-of-world local currency price index rose by less than 50 percent.


% Source PDF 11, text block 2.
The previous figure suggests a minor role for exchange rate movements in driving the recent collapse in the US net foreign asset position. However, US investors buying foreign equity do not necessarily hold the MSCI ROW index. To more precisely assess the role of exchange rate movements on net equity revaluations, we turn to Table 1.3 of the “International Investment Position” data (Bureau of Economic Analysis 2024), which breaks down valuation changes in all components of US foreign assets and liabilities into a part accounted for by exchange rate changes and a part that reflects local currency price movements. In Table 1 we use these data to decompose changes in equity valuations into these two components for our two subperiods.\hyperlink{fn:14}{\textsuperscript{14}}


% Source PDF 11, text block 3.
In the first subperiod, the US dollar was depreciating, and dollar depreciation amplified the dollar value of US equity abroad. This exchange rate–driven revaluation improved the US net foreign asset position by 15.8 percent of US corporate gross value added, out of a total revaluation of the net equity position of 40.4 percent of GVA. Thus, confirming the impression from Figure 6, exchange rate movements accounted for a large part of total equity revaluation.


% Source PDF 11, text block 4.
Between 2010 and 2020, dollar appreciation led to a decline in the dollar value of US foreign equity assets cumulating to 10.1 percent of corporate GVA. However, the dominant driver of net equity revaluation in this period was the sharp increase in the value of US equity liabilities, cumulating to 103.6 percent of corporate GVA, for which exchange rate movements obviously play no role. Thus, dollar appreciation accounted for only about one-sixth of the total net decline in the US net equity position over this period.


% Source PDF 11, text block 5.
Measurement of Cross-Border Positions and Income Flows.—We have relied on data from the Integrated Macroeconomic Accounts (IMA) to measure cross-border asset positions on a residence basis. However, all estimates of cross-border asset


% Source PDF 11, text block 6.
{\footnotesize \hypertarget{fn:14}{}\textbf{Footnote 14.} For each subperiod, we cumulate the revaluation of US equity assets and US equity liabilities that are accounted by exchange rate movements and price changes and divide them by US corporate gross value added at the end of the period.\par}



% ===== ORIGINAL PDF PAGE 12; JOURNAL PAGE 2162 =====
\sourcepage{12}{2162}

% Source PDF 12, text block 1.
positions, including those in the IMA, are subject to a range of measurement challenges. Many of these arise due to the difficulties in dividing up the activities of multinational corporations into components attributable to their subsidiaries resident in different countries. Bertaut, Bressler, and Curcuru (2019) provides an excellent overview of the difficulties in measuring corporate activity on a residence basis and in measuring cross-border equity positions.\hyperlink{fn:15}{\textsuperscript{15}} Households shelter wealth in low tax offshore jurisdictions (Zucman 2013), and firms use subsidiaries in the same places to raise capital (Coppola et al. 2021), muddying the official map of cross-border positions. Another issue specific to closely held foreign direct investment equity is uncertainty over appropriate market valuation.


% Source PDF 12, text block 2.
A separate but related set of measurement concerns has to do with the fact that while the US net foreign asset position is large and negative, US primary income from abroad as measured in the current account remains positive (see, e.g., Curcuru, Thomas, and Warnock 2013 and Guvenen et al. 2022).


% Source PDF 12, text block 3.
In Section 1 of the Supplemental Appendix, we provide more discussion of the literature on these measurement issues, and we explore the sensitivity of our model exercises to smaller estimates of gross cross-border equity positions using an adjustment based on Bertaut, Bressler, and Curcuru (2019).


% Source PDF 12, text block 4.
Flows, Stocks, and Valuation of the US Corporate Sector.—We have shown that changes in US equity values play a very important role, in an accounting sense, in explaining changes in the US NFA position. We will build an economic model to interpret the causes and welfare consequences of these valuation effects. That model will incorporate a range of additional data that we now describe. We first define measurement concepts that are consistent across the model and data for flows including value added, taxes, labor compensation, investment, earnings, and free cash flow. We then discuss our measurement of stocks, including the reproduction value of the stock of capital in the corporate sector, and the value of corporations. Our primary measure of income to firm owners, which we label “free cash flow,” is the natural one in the context of our model: value added minus the sum of labor compensation, investment, and taxes paid. Our primary valuation measure is the aggregate value of the corporate sector’s nonfinancial assets that generate this income, which we label “enterprise value.” We will interpret these broad empirical measures for corporate income and valuation through the lens of a model in which firms are 100 percent equity financed. In doing so, we are implicitly assuming that corporate valuation is independent of how firms are financed.\hyperlink{fn:16}{\textsuperscript{16}}


% Source PDF 12, text block 5.
Flows in the US Corporate Sector.—We use Tables S.5 and S.6 of the Integrated Macroeconomic Accounts to measure the flows and balance sheets of the US corporate sector. Table S.5 presents data for the nonfinancial corporate business sector, and Table S.6 presents data for the financial business sector. The overwhelming portion of foreign portfolio and direct investment into the United States is directed toward these two sectors. We combine these two accounts into an aggregated corporate sector. The national accounts follow the residence principle. Thus, the value


% Source PDF 12, text block 6.
{\footnotesize \hypertarget{fn:15}{}\textbf{Footnote 15.} See also Avdjiev et al. (2018) and Lane (2020).\par}

{\footnotesize\hypertarget{fn:16}{}\textbf{Footnote 16.} Introducing corporate debt finance in our model would not impact model enterprise value because our model features no frictions that would break Modigliani-Miller.\par}



% ===== ORIGINAL PDF PAGE 13; JOURNAL PAGE 2163 =====
\sourcepage{13}{2163}

% Source PDF 13, text block 1.
added of US-resident affiliates of foreign multinationals is counted as part of US value added, while the value added by US-owned businesses abroad is not.


% Source PDF 13, text block 2.
The gross value added of these sectors is divided into four categories of income in Tables S.5 and S.6: consumption of fixed capital (depreciation), compensation of employees, taxes on production and imports less subsidies, and net operating surplus. We measure the earnings of the corporate sector as net operating surplus less current taxes on income and wealth. We measure the free cash flow of the corporate sector as net operating surplus less current taxes on income and wealth less net capital formation. Free cash flow corresponds to the after-tax cash flow from operations of corporations resident in the United States that is available to be paid out to investors in the debt and equity of those corporations. In the data, only some of this cash flow is actually paid out to investors, while the rest of it is used to acquire, on net, financial assets (as accounted for in Tables S.5 and S.6). Thus, our empirical measure of free cash flow corresponds to what dividends would be if firms were 100 percent equity financed and maintained no financial assets.


% Source PDF 13, text block 3.
In Figure 7, we examine the ratio of our measure of free cash flow to US corporate sector GVA. We see that this ratio has risen substantially over the past 14 years, compared with the period before 2007. This increase in payouts arises from a combination of reductions in taxes, labor compensation, and investment as ratios to corporate gross value added. The path for earnings reported by Greenwald, Lettau, and Ludvigson (2025) in their Figure 4b looks broadly similar to our series for free cash flow. But note that their earnings measure subtracts depreciation from gross operating surplus, while our cash flow concept subtracts investment, which they do not model. In Figure A4 in Appendix A.A3, we document a similar increase in free cash flow when restricting attention to the nonfinancial corporate sector. In common with Greenwald, Lettau, and Ludvigson (2025) and Atkeson, Heathcote, and Perri (2025), we will argue that rising US earnings are central to explaining rising US corporate valuations. Our novel argument, relative to those papers, is that this rise in valuations also drives the US NFA position.


% Source PDF 13, text block 4.
Valuation and Capital in the US Corporate Sector.—We now describe how we measure the enterprise value of the nonfinancial assets held by US resident corporations, where these assets generate the free cash flow described above. To measure enterprise value, we make several adjustments to the balance sheet data for the corporate sector presented in Tables S.5 and S.6 of the IMA. The following stylized balance sheet for the US corporate sector is useful for organizing our discussion of these adjustments.\hyperlink{fn:17}{\textsuperscript{17}}


% Source PDF 13, text block 5.
\begin{center}\small
\textbf{Corporate Sector Balance Sheet}\\[4pt]
\begin{tabular}{p{0.43\linewidth}|p{0.43\linewidth}}\toprule
\centering Assets & \centering Liabilities\tabularnewline\midrule
Nonfinancial assets (replacement or enterprise value) & Equity (measured at market value)\\[4pt]
Financial assets (includes US FDI in ROW) & Financial liabilities (debt, bank loans, etc., including ROW FDI in US)\\\bottomrule
\end{tabular}
\end{center}


% Source PDF 13, text block 9.
{\footnotesize \hypertarget{fn:17}{}\textbf{Footnote 17.} Recall that this balance sheet is an aggregate of both US firms with overseas subsidiaries (i.e., the parent firm is in the United States) and US-resident subsidiaries of foreign multinationals.\par}



% ===== ORIGINAL PDF PAGE 14; JOURNAL PAGE 2164 =====
\sourcepage{14}{2164}

% Figure 7: original PDF page 14, plot crop x=60..444, y=62..290 points.
\begin{center}
\begingroup\renewcommand{\thefigure}{7}\refstepcounter{figure}\label{fig:7}\endgroup
\includegraphics[page=14,trim=60 430 60 62,clip,width=0.95\linewidth]{\sourcepdf}
\par\small Figure 7. US Corporate Free Cash Flow to Owners
\end{center}


% Source PDF 14, text block 1.
Our specific aim is to value the nonfinancial assets held by US resident corporations, corresponding to the first entry in the left column of this balance sheet. We consider two measures of this value. The first of these is a measure of the replacement value of these nonfinancial assets. This measure is reported directly in the Integrated Macroeconomic Accounts.


% Source PDF 14, text block 2.
The second enterprise value measure captures the value that financial markets attach to corporate nonfinancial assets located in the United States. It is measured as the sum of the market value of resident corporations’ equities plus the value of their financial liabilities (both on the right side of the corporate sector balance sheet) less the value of financial assets on the left side of this balance sheet.\hyperlink{fn:18}{\textsuperscript{18}}’\hyperlink{fn:19}{\textsuperscript{19}}


% Source PDF 14, text block 3.
The financial assets of these firms, listed as the second entry on the left side of this balance sheet, include the usual financial instruments as well as the debt and equity components of US parent firms’ foreign direct investment abroad. The financial liabilities of these firms, listed as the second item on the right side of this balance sheet, include the usual financial instruments plus the debt and equity components of the direct investment of foreign parent firms into their US subsidiaries. Excluding US FDI in the rest of the world from enterprise value but including rest-of-world FDI into the United States aligns our measure of US enterprise value with the residence principle.


% Source PDF 14, text block 4.
{\footnotesize \hypertarget{fn:18}{}\textbf{Footnote 18.} Note that it would not be appropriate to equate the value of US-located firms to the value of equity alone because some fraction of firms’ future cash flow is pledged to debt holders and bank lenders, and thus, some fraction of firm value is reflected in the value of those liabilities.\par}


% Source PDF 14, text block 5.
{\footnotesize \hypertarget{fn:19}{}\textbf{Footnote 19.} Our enterprise value concept is roughly similar to the concept of enterprise value used as a valuation benchmark for individual companies and is closely related to that used in Hall (2001). It is also similar to the measures of the market values of the nonfinancial assets of US nonfinancial and financial corporations presented in Table B1 of the Financial Accounts of the United States.\par}



% ===== ORIGINAL PDF PAGE 15; JOURNAL PAGE 2165 =====
\sourcepage{15}{2165}

% Figure 8: original PDF page 15, plot crop x=60..444, y=62..228 points.
\begin{center}
\begingroup\renewcommand{\thefigure}{8}\refstepcounter{figure}\label{fig:8}\endgroup
\includegraphics[page=15,trim=60 492 60 62,clip,width=0.95\linewidth]{\sourcepdf}
\par\small Figure 8. Values and Cash Flow, US Corporate Sector
\end{center}


% Source PDF 15, text block 1.
In the left panel of Figure 8, we show the ratio of enterprise value to value added for the US corporate sector (blue line) and the ratio of the replacement value of the stock of capital in those corporations to value added (red line). The figure indicates that the capital-output ratio has been quite stable over time, while the enterprise value of US corporations has risen dramatically. A direct implication of the divergence between these two lines is that Tobin’s Q for the US corporate sector, measured as the ratio of enterprise value to replacement value of the capital stock, has risen substantially over the 2010–2023 period. Figure A3 in Appendix A.A3 plots the same series for the nonfinancial corporate sector and shows that the dynamics for valuations are very similar to those for the entire corporate sector. Figure A5 in Appendix A.A3 shows that the time path for the market value of equity for the nonfinancial corporate sector tracks enterprise value closely but is lower by around 50 percent of GDP, indicating that corporate financial liabilities exceed financial assets (see the stylized balance sheet above).


% Source PDF 15, text block 2.
In the right panel of Figure 8, we construct a measure of the payout yield for the US corporate sector based on the ratio of the free cash flow to firm owners to our measure of the enterprise value of this sector. What is striking about this figure is that the ratio of free cash flow to enterprise value has not changed much for most of the past decade, relative to the period before 2007. Thus, it appears that a substantial portion of the increase in the ratio of the enterprise value of US corporations to GVA can be accounted for by an increase in the ratio of free cash flow to GVA.


% Source PDF 15, text block 3.
We now document that the sharp increases in corporate enterprise value and payouts just described are US-specific phenomena and that similar increases have not occurred in the rest of the developed world. In the left panel of Figure 9, we plot, for the period 1990–2022, corporate sector enterprise values and free cash flows for the United States, for an aggregate of the other G6 countries, and for the European Union. The figure highlights the divergence between the ratios in the United States and in the rest of the world. Since the Great Recession, the ratio of enterprise value to GVA in the United States rose from 2 to over 4.5, while the ratio abroad has been essentially constant over the same time span. The right panel of Figure 9 shows that the diverging patterns in corporate sector enterprise value are mirrored by diverging



% ===== ORIGINAL PDF PAGE 16; JOURNAL PAGE 2166 =====
\sourcepage{16}{2166}

% Figure 9: original PDF page 16, plot crop x=60..444, y=62..225 points.
\begin{center}
\begingroup\renewcommand{\thefigure}{9}\refstepcounter{figure}\label{fig:9}\endgroup
\includegraphics[page=16,trim=60 495 60 62,clip,width=0.95\linewidth]{\sourcepdf}
\par\small Figure 9. Values and Cash Flow, United States, G6, and European Union
\end{center}


% Source PDF 16, text block 1.
paths in free cash flow. From 2007 to 2022, the ratio of free cash flow to GVA in the United States nearly doubled, while abroad, this ratio over the same period was largely unchanged. A natural interpretation is that differential dynamics of free cash flow are key to explaining the differential behavior of United States and foreign equity markets.\hyperlink{fn:20}{\textsuperscript{20}}


% Source PDF 16, text block 2.
\section*{II. Model}
\addcontentsline{toc}{section}{II. Model}


% Source PDF 16, text block 3.
We now develop a simple international macrofinance model of flows, stocks, and valuations of the US corporate sector and of the US current account and NFA position. We use this model to measure the factors driving observed flows, stocks, and valuations of the US corporate sector, the US current account, and the US NFA position over the period from 1990 through the third quarter of 2023. Then, we conduct counterfactual exercises relative to this model baseline to consider how these driving factors impacted the welfare of US households.


% Source PDF 16, text block 4.
The model has two regions: a domestic economy we think of as the United States and a foreign economy that stands in for the rest of the world. Each region is populated by a continuum of identical households. Heterogeneous firms in each economy produce a continuum of nontradable intermediate varieties. These intermediates are combined to produce a single composite final good that is traded internationally and used for consumption and investment. Intermediates-producing firms enjoy pricing power and hence earn factorless income. Households receive labor income, dividends from holdings of corporate equity both in their region of residence and abroad, and interest income from a risk-free bond that is traded internationally.


% Source PDF 16, text block 5.
{\footnotesize \hypertarget{fn:20}{}\textbf{Footnote 20.} Details on how we constructed data in Figure 9 are in Appendix A.A5. This evidence that outside the United States, corporate free cash flow has not risen relative to GVA is consistent with other studies. See, for example, Lequiller and Blades (2014, chap. 3); Philippon (2019, chap. 5); and Gutiérrez and Philippon (2023). Gutiérrez and Piton (2020) find evidence of labor’s share declining much more in the United States than in other advanced economies. In contrast, De Loecker and Eeckhout (2021), using the Worldscope dataset of firm financial statements, argue that markups and profits have risen in Europe as well as in the United States.\par}



% ===== ORIGINAL PDF PAGE 17; JOURNAL PAGE 2167 =====
\sourcepage{17}{2167}

% Source PDF 17, text block 1.
In our baseline model specification, we assume that foreign households are risk neutral and that their rate of time preference determines the cost of capital for firms worldwide. That assumption allows us to characterize equilibrium allocations in closed form and to illustrate the economic mechanisms at work as transparently as possible. We also assume that both countries produce and consume the same final good, so the terms of trade and the real exchange rate in the model will always be equal to one. Recall that exchange rate movements account for only a small portion of the valuation effects in the NFA position between 2010 and 2022. In Section 5 of the Supplemental Appendix, we discuss a generalization of the model in which foreign- and domestically produced goods are imperfect substitutes, in which case shocks to monopoly power and/or productivity have the potential to affect the terms of trade.


% Source PDF 17, text block 2.
\subsection*{A. Intermediate-Goods Firms}
\addcontentsline{toc}{subsection}{A. Intermediate-Goods Firms}


% Source PDF 17, text block 3.
In each country there is a unit mass of different intermediate varieties indexed by i  $\in [0,1]$. Let $Y_{it}$ denote total production of variety $i$ at date $t$. Domestic output of the final good is given by


% Source PDF 17, text block 4.
\setcounter{equation}{3}
\begin{equation}\label{eq:4}
Y_t=\left(\int_0^1Y_{it}^{\left(\frac{\varepsilon-1}{\varepsilon}\right)}\,di\right)^{\frac{\varepsilon}{\varepsilon-1}},
\end{equation}


% Source PDF 17, text block 9.
where $\varepsilon >1$ is the elasticity of substitution in production between different varieties.


% Source PDF 17, text block 10.
Output can be consumed domestically, exported, or transformed into investment. Thus,


% Source PDF 17, text block 11.
\[
C_t+G_t+Q_tX_t+NX_t=Y_t,
\]


% Source PDF 17, text block 12.
where $C_{t}$ is private consumption, $G_{t}$ is public consumption, $Q_{t}X_{t}$ is investment expenditure in units of final output, and $NX_{t}$ is net exports. Investment goods augment the capital stock in the standard way:


% Source PDF 17, text block 13.
\[
K_{t+1}=(1-\delta_t)K_t+X_t,
\]


% Source PDF 17, text block 14.
where $\delta _{t}$ is a time-varying depreciation rate. The replacement value of the capital stock at the end of period $t$ in units of final consumption is denoted by $Q_{t}K_{t+1}$, where this replacement value evolves over time according to


% Source PDF 17, text block 15.
\setcounter{equation}{4}
\begin{equation}\label{eq:5}
Q_tK_{t+1}=Q_{t-1}K_t+(Q_t-Q_{t-1})K_t-\delta_t Q_tK_t+Q_tX_t,
\end{equation}


% Source PDF 17, text block 16.
where $Q_{t-1}K_{t}$ is the replacement value of the capital stock at the end of period $t-1,(Q_{t}-Q_{t-1})K_{t}$ is revaluation of installed capital between $t-1$ and $t$, and the term $\delta _{t}Q_{t}K_{t}$ corresponds to consumption of fixed capital in units of the final consumption good at $t$.


% Source PDF 17, text block 17.
Both countries produce the same final good. Thus, if we use asterisks to denote foreign variables, the world resource constraint is


% Source PDF 17, text block 18.
\[
C_t+C_t^*+G_t+G_t^*+Q_tX_t+Q_t^*X_t^*=Y_t+Y_t^*.
\]



% ===== ORIGINAL PDF PAGE 18; JOURNAL PAGE 2168 =====
\sourcepage{18}{2168}

% Source PDF 18, text block 1.
Within each country, there are two sorts of firms that can produce a given variety of intermediate good: a single leader firm with productivity $z_{Ht}$, and a fringe of identical follower firms, each with productivity $z_{Lt}\leq z_{Ht}$. An intermediate firm with productivity $z_{t}$ that rents capital $k_{t}$ and labor $l_{t}$ produces output $y_{t}$, given by


% Source PDF 18, text block 2.
\[
y_t=z_tk_t^{\alpha_t}(Z_tl_t)^{1-\alpha_t},
\]


% Source PDF 18, text block 3.
where $Z_{t}$ is aggregate labor productivity common to all firms and where $\alpha _{t}$ is a time-varying parameter determining the relative importance of capital versus labor in production costs. The production technologies for leader and follower firms are common across all intermediate varieties in the United States. We denote the corresponding productivities in the ROW by $z_{Ht}^{*},z_{Lt}^{*}$.

Bertrand price competition between the leader firm and the follower firms for each variety determines the markup of price over marginal cost charged by the leader firm, as in Bernard et al. (2003); Atkeson and Burstein (2007); and Peters (2020). Specifically, let $R_{t}$ and $W_{t}$ denote the domestic rental rates for capital and labor in the United States. These factor prices, together with firm productivity $z_{t}$, determine intermediate firms’ unit cost. Intermediate firms pay a proportional tax at rate $\tau _{t}$ on sales. Because these firms have no intermediate inputs, this can be interpreted as a value-added tax. The leader firms for each variety move first and set a price $p_{it}$. If these firms did not face any latent competition from follower firms, they would solve the standard monopolistic competition profit-maximization problem, and the after-tax price would be a markup over marginal (and average) cost of $\varepsilon /(\varepsilon -1)$. However, the leader firm also recognizes that if it sets its price above the marginal cost of the firm with productivity $z_{Lt}$, then latent competitors will be able to profitably enter and will in fact corner the market. Thus, the leader firm effectively faces an additional constraint on pricing, one that ensures that competitors do not enter and the leader retains a 100 percent market share. Given these constraints on pricing, the equilibrium output wedge $\mu _{t}$ between after-tax revenues relative to total costs is given by


% Source PDF 18, text block 4.
\setcounter{equation}{5}
\begin{equation}\label{eq:6}
\mu_t=\frac{(1-\tau_t)p_{it}}{\mathit{cost}_t(z_{Ht})}=\min\left\{\frac{\varepsilon}{\varepsilon-1},\frac{\mathit{cost}_t(z_{Lt})}{\mathit{cost}_t(z_{Ht})}=\frac{z_{Ht}}{z_{Lt}}\right\},
\end{equation}


% Source PDF 18, text block 8.
where $\mathit{cost}_{t}(z_{Ht})$ is the unit cost of production at $t$ for a leader firm. We assume that $z_{Ht}/z_{Lt}<\varepsilon /(\varepsilon -1)$ for all $t$ so that the output wedge is always driven by the threat of potential competition, $\mu _{t}=z_{Ht}/z_{Lt}$.


% Source PDF 18, text block 9.
Note that because all varieties are symmetric, equilibrium prices, output wedges, labor, capital, and output are identical across varieties: $p_{it}=P_{t},k_{it}=K_{t},l_{it}=L_{t}$, and


% Source PDF 18, text block 10.
\setcounter{equation}{6}
\begin{equation}\label{eq:7}
y_{it}=Y_{it}=Y_t=z_{Ht}K_t^{\alpha_t}(Z_tL_t)^{1-\alpha_t}.
\end{equation}


% Source PDF 18, text block 11.
Without loss of generality, we normalize $P_{t}=1$ for all $t$. In our baseline model, we also assume exogenous and fixed labor supply and normalize $L_{t}=1$.


% Source PDF 18, text block 12.
Tax payments by intermediate goods firms fund government purchases: $G_{t}=\tau _{t}Y_{t}$.



% ===== ORIGINAL PDF PAGE 19; JOURNAL PAGE 2169 =====
\sourcepage{19}{2169}

% Source PDF 19, text block 1.
Since the production function for final output in equation (4) has constant returns to scale, in equilibrium, final output is equal to the pretax revenue of intermediate goods firms. After-tax revenue from intermediate firms is divided between wage payments to labor, rental payments to capital, and factorless income, which we denote by $\Pi _{t}$. The share of pretax output accruing as factorless income to owners of intermediate goods firms is


% Source PDF 19, text block 2.
\setcounter{equation}{7}
\begin{equation}\label{eq:8}
\frac{\Pi_t}{Y_t}=\left(\frac{\mu_t-1}{\mu_t}\right)(1-\tau_t),
\end{equation}


% Source PDF 19, text block 4.
while the shares going to labor and capital are


% Source PDF 19, text block 5.
\setcounter{equation}{8}
\begin{equation}\label{eq:9}
\frac{W_tL_t}{Y_t}=\frac{(1-\alpha_t)}{\mu_t}(1-\tau_t),
\end{equation}


% Source PDF 19, text block 6.
\setcounter{equation}{9}
\begin{equation}\label{eq:10}
\frac{R_tK_t}{Y_t}=\frac{\alpha_t}{\mu_t}(1-\tau_t),
\end{equation}


% Source PDF 19, text block 7.
and the remaining share $\tau _{t}$ goes to taxes.


% Source PDF 19, text block 8.
\subsection*{B. Investment Firms}
\addcontentsline{toc}{subsection}{B. Investment Firms}


% Source PDF 19, text block 9.
In addition to intermediates-producing firms, a second set of competitive firms holds and rents out capital and makes investment choices. These competitive capital-managing firms choose investment to maximize the expected present value of their dividends. Dividends from these firms are given by


% Source PDF 19, text block 10.
\setcounter{equation}{10}
\begin{equation}\label{eq:11}
D_t^K=R_tK_t-Q_tX_t=R_tK_t-Q_tK_{t+1}+Q_t(1-\delta_t)K_t.
\end{equation}


% Source PDF 19, text block 11.
Investment firms discount cash flow one period ahead at rate $r_{t+1}^{*}$. At each date $t$, given $K_{t}$, they choose $K_{t+1}$ to solve


% Source PDF 19, text block 12.
\[
\max_{K_{t+1}}\left\{-Q_tK_{t+1}+\frac{1}{1+r_{t+1}^*}\mathbb{E}_t\left[R_{t+1}K_{t+1}+(1-\delta_{t+1})Q_{t+1}K_{t+1}\right]\right\},
\]


% Source PDF 19, text block 15.
where the interpretation is that purchasing one more unit of new capital at $t$ reduces current dividends by the price of capital $Q_{t}$ but generates additional rental income $R_{t+1}$ and a resale value of undepreciated capital $(1-\delta _{t+1})Q_{t+1}$ in the next period.


% Source PDF 19, text block 16.
The first-order condition to this problem is


% Source PDF 19, text block 17.
\setcounter{equation}{11}
\begin{equation}\label{eq:12}
Q_t=\frac{1}{1+r_{t+1}^*}\mathbb{E}_t\left[R_{t+1}+(1-\delta_{t+1})Q_{t+1}\right].
\end{equation}


% Source PDF 19, text block 19.
In our model, we assume that all firms are financed entirely by equity and have no financial assets. Thus, the measure of aggregate dividends paid by US firms in the model corresponds to a measure of free cash flow from operations available to be paid to all investors in the firm,


% Source PDF 19, text block 20.
\setcounter{equation}{12}
\begin{equation}\label{eq:13}
D_t=\Pi_t+D_t^K,
\end{equation}



% ===== ORIGINAL PDF PAGE 20; JOURNAL PAGE 2170 =====
\sourcepage{20}{2170}

% Source PDF 20, text block 1.
and likewise for foreign dividends. The measure of firm value $V_{t}$ in the model corresponds to the market valuation of these free cash flows from operations.


% Source PDF 20, text block 2.
We refer to the after-tax net operating surplus of firms in our model as the earnings of these firms. These earnings are given by


% Source PDF 20, text block 3.
\setcounter{equation}{13}
\begin{equation}\label{eq:14}
E_t=(1-\tau_t)Y_t-W_tK_t-\delta_tQ_tK_t.
\end{equation}
% Conversion note: equation (14) prints W_t K_t, although the surrounding wage-bill notation uses W_t L_t. Preserved as printed.



% Source PDF 20, text block 4.
Note that our measure of aggregate dividends $D_{t}$ is equal to our measure of earnings $E_{t}$ less net investment $Q_{t}X_{t}-\delta _{t}Q_{t}K_{t}$, as is standard.


% Source PDF 20, text block 5.
The profit-maximization problems for foreign firms mirror those for domestic ones. Foreign technology parameters are mostly identical to domestic ones, with the exceptions of the intermediate firm productivity values, $z_{Ht}^{*}$ and $z_{Lt}^{*}$ (and thus the output wedge $\mu ^{*}=z_{Ht}^{*}/z_{Lt}^{*}$), and the replacement cost of capital, $Q_{t}^{*}$.


% Source PDF 20, text block 6.
\subsection*{C. Households}
\addcontentsline{toc}{subsection}{C. Households}


% Source PDF 20, text block 7.
Lifetime utility for the domestic representative infinitely lived household is given by


% Source PDF 20, text block 8.
\setcounter{equation}{14}
\begin{equation}\label{eq:15}
\mathbb{E}_0\sum_{t=0}^{\infty}\left(\frac{1}{1+\rho}\right)^t\log(C_t),
\end{equation}


% Source PDF 20, text block 13.
where $\rho $ is the constant rate of time preference.


% Source PDF 20, text block 14.
The assets in this economy are shares in domestic and foreign firms and a one-period nominal bond. Domestic households enter period $t$ owning a fraction $\lambda _{t-1}$ of shares in domestic firms (foreign households own fraction $1-\lambda _{t-1}$) and a fraction $\lambda _{t-1}^{*}$ of foreign firms. They also enter period $t$ with $B_{t}$ units of bonds, which pay interest at rate $r_{t}^{*}$. Each period, households buy new domestic and foreign shares at prices $V_{t}$ and $V_{t}^{*}$, and bonds $B_{t+1}$ at a price normalized to one. The interest rate between $t$ and $t+1,r_{t+1}^{*}$, is known at date $t$ (bonds are risk free) and is the same rate used by firms to discount future cash flows. The flow budget constraint for the domestic representative household is


% Source PDF 20, text block 15.
\setcounter{equation}{15}
\begin{equation}\label{eq:16}
\begin{aligned}
C_t+(\lambda_t-\lambda_{t-1})V_t+(\lambda_t^*-\lambda_{t-1}^*)V_t^*+B_{t+1}
&=W_tL_t+\lambda_{t-1}D_t\\
&\quad+\lambda_{t-1}^*D_t^*+(1+r_t^*)B_t.
\end{aligned}
\end{equation}


% Source PDF 20, text block 17.
Foreign households are symmetric to domestic ones, except that we assume they have linear utility $(u^{*}(C_{t}^{*})=C_{t}^{*})$ and a time-varying discount factor $\rho _{t}^{*}$. The foreign discount factor between $t$ and $t+1,\rho _{t+1}^{*}$, is known at $t$.

Note that because foreign households have linear utility, the world interest rate is pinned down at


% Source PDF 20, text block 18.
\setcounter{equation}{16}
\begin{equation}\label{eq:17}
r_{t+1}^*=\rho_{t+1}^*
\end{equation}


% Source PDF 20, text block 19.
for all dates $t$.



% ===== ORIGINAL PDF PAGE 21; JOURNAL PAGE 2171 =====
\sourcepage{21}{2171}

% Source PDF 21, text block 1.
\subsection*{D. Expectations}
\addcontentsline{toc}{subsection}{D. Expectations}


% Source PDF 21, text block 2.
Equilibrium investment in capital and the valuation of the future streams of dividends for US and foreign firms depend on expectations of future realizations of the parameters of the model. We assume that households and firms make decisions taking as given forecasts for the evolution of all model parameter values very similar to those assumed in the transition experiment in Farhi and Gourio (2018). We assume households perceive no uncertainty around these forecasts.


% Source PDF 21, text block 3.
We now describe these forecasts. At each date $t$, model agents:


% Source PDF 21, text block 4.
(i) Observe the one-period-ahead discount rate, $r_{t+1}^{*}=\rho _{t+1}^{*}$. They expect $\mathbb{E}_{t}[r_{t+j}^{*}]=r_{t+1}^{*}$ for all $j\geq 2$.


% Source PDF 21, text block 5.
(ii) Perfectly forecast one-period-ahead growth in global labor productivity, $g_{t+1}$.


% Source PDF 21, text block 6.
Thus, $\mathbb{E}_{t}[Z_{t+1}/Z_{t}]=Z_{t+1}/Z_{t}=1+g_{t+1}$.


% Source PDF 21, text block 7.
(iii) Expect global productivity to grow at a constant rate from period $t+1$


% Source PDF 21, text block 8.
onward. We denote the future expected trend growth rate at $t$ by $\bar{g}_{t+1}$. Thus, $\mathbb{E}_{t}[Z_{t+j}/Z_{t+j-1}]=1+\bar{g}_{t+1}$ for all $j\geq 2$. Note that we do not impose $\bar{g}_{t+1}=g_{t+1}$.


% Source PDF 21, text block 9.
(iv) Perfectly forecast next-period values for leader and follower productivities


% Source PDF 21, text block 10.
in the two economies, $z_{H,t+1}$ and $z_{L,t+1}$ and $z_{H,t+1}^{*}$ and $z_{L,t+1}^{*}$. In addition, we assume that they expect these values to persist: $\mathbb{E}_{t}[z_{H,t+j}]=z_{H,t+1},\mathbb{E}_{t}[z_{L,t+j}]=z_{L,t+1},\mathbb{E}_{t}[z_{H,t+j}^{*}]=z_{H,t+1}^{*}$, and $\mathbb{E}_{t}[z_{L,t+j}^{*}]=z_{L,t+1}^{*}$ for all $j\geq 2$. Thus, agents expect constant output wedges, $\mu _{t+1}$ and $\mu _{t+1}^{*}$, from period $t+1$ on.


% Source PDF 21, text block 11.
(v) Perfectly forecast next-period values for the technological parameters $\alpha _{t+1}$


% Source PDF 21, text block 12.
and $\delta _{t+1}$ and the tax rate $\tau _{t+1}$. In addition, they expect these parameter values to persist: $\mathbb{E}_{t}[\alpha _{t+j}]=\alpha _{t+1},\mathbb{E}_{t}[\delta _{t+j}]=\delta _{t+1}$, and $\mathbb{E}_{t}[\tau _{t+j}]=\tau _{t+1}$ for all j  $\geq 2$.


% Source PDF 21, text block 13.
(vi) Expect no changes in the relative prices of investment goods: $\mathbb{E}_{t}[Q_{t+j}]=$


% Source PDF 21, text block 14.
$Q_{t}$ and $\mathbb{E}_{t}[Q_{t+j}^{*}]=Q_{t}^{*}$ for all $j\geq 1$.


% Source PDF 21, text block 15.
To summarize, in each period $t$, agents receive news about the cost of capital $r_{t+1}^{*}$ at $t$ and values of other model parameters that will be realized at $t+1$. They treat these parameters as if they followed a random walk. That is, their expectations for the values of these parameters at dates $t+j$ are equal to the value that they expect at $t+1$.


% Source PDF 21, text block 16.
\subsection*{E. Asset Pricing}
\addcontentsline{toc}{subsection}{E. Asset Pricing}


% Source PDF 21, text block 17.
Firm value $V_{t}$ in the model can be decomposed into the ex-dividend value of claims to investment-producing firms, plus the value of claims to profits from intermediate-goods firms. As is standard in a model with constant returns to scale and no investment adjustment costs, equation (12) implies that the present value at $t$ of dividends from capital-managing firms from $t+1$ on is equal to the expected replacement value of capital in the next period. That is, $V_{t}^{K}=Q_{t}K_{t+1}$.



% ===== ORIGINAL PDF PAGE 22; JOURNAL PAGE 2172 =====
\sourcepage{22}{2172}

% Source PDF 22, text block 1.
The ex-dividend price of a share of all domestic intermediate-good-producing firms is the expected present value of the future stream of monopoly profits these firms will earn. Given our assumptions, agents know at $t$ the parameters determining profits at $t+1,\Pi _{t+1}$. They expect the share of income corresponding to after-tax profits to remain constant from $t+1$ onward, and they expect income (and thus profits) to grow at constant rate $\bar{g}_{t+1}$. They discount future profit income at a constant rate $r_{t+1}^{*}$. Thus,


% Source PDF 22, text block 2.
\setcounter{equation}{17}
\begin{equation}\label{eq:18}
V_t^{\Pi}=\mathbb{E}_t\left[\sum_{j=1}^{\infty}\frac{\Pi_{t+j}}{(1+r_{t+1}^*)^j}\right]=\frac{\Pi_{t+1}}{r_{t+1}^*-\bar g_{t+1}}.
\end{equation}


% Source PDF 22, text block 8.
Thus, the market price of all domestic firms is given by\hyperlink{fn:21}{\textsuperscript{21}}


% Source PDF 22, text block 9.
\setcounter{equation}{18}
\begin{equation}\label{eq:19}
V_t=V_t^K+V_t^{\Pi}=Q_tK_{t+1}+\frac{\Pi_{t+1}}{r_{t+1}^*-\bar g_{t+1}}.
\end{equation}


% Source PDF 22, text block 11.
\subsection*{F. Balance of Payments Accounting}
\addcontentsline{toc}{subsection}{F. Balance of Payments Accounting}


% Source PDF 22, text block 12.
We now consider our model’s implications for the current account and the US NFA position. The current account is defined as national savings minus investment, where national savings is comprised of household saving plus government saving plus corporate saving. In our model, the government runs a balanced budget period by period, and corporate saving is identical to corporate investment. Thus, the model current account is identical to household saving. The change in the NFA position of the United States in the model is the sum of the current account and the revaluations of cross-border asset holdings of households in the United States and in the ROW.


% Source PDF 22, text block 13.
We now solve for the savings of US households. The representative US household at each date $t$ chooses a sequence for consumption to maximize utility (15), subject to a lifetime budget constraint. Given that this household has logarithmic utility, it consumes a constant fraction of its lifetime wealth inclusive of current income. That is, at each date $t$,


% Source PDF 22, text block 14.
\setcounter{equation}{19}
\begin{equation}\label{eq:20}
C_t=\frac{\rho}{1+\rho}\mathit{Wealth}_t,
\end{equation}


% Source PDF 22, text block 15.
where household wealth at $t$ inclusive of current income is given by


% Source PDF 22, text block 16.
\setcounter{equation}{20}
\begin{equation}\label{eq:21}
\begin{aligned}
\mathit{Wealth}_t={}&W_tL_t+H_t+\lambda_{t-1}D_t+\lambda_{t-1}^*D_t^*\\
&+\lambda_{t-1}V_t+\lambda_{t-1}^*V_t^*+(1+r_t^*)B_t.
\end{aligned}
\end{equation}


% Source PDF 22, text block 17.
Here, $H_{t}$ denotes human wealth excluding current labor earnings and is given by


% Source PDF 22, text block 18.
\setcounter{equation}{21}
\begin{equation}\label{eq:22}
H_t\equiv\frac{W_{t+1}L_{t+1}}{r_{t+1}^*-\bar g_{t+1}}.
\end{equation}


% Source PDF 22, text block 20.
{\footnotesize \hypertarget{fn:21}{}\textbf{Footnote 21.} In our model, firms issue no debt, and thus, $V_{t}$ is both the market value of equity and enterprise value. If we were to introduce corporate debt, equation (19) would still be the expression for enterprise value.\par}



% ===== ORIGINAL PDF PAGE 23; JOURNAL PAGE 2173 =====
\sourcepage{23}{2173}

% Source PDF 23, text block 1.
US household saving is the difference between households’ current income and consumption. Using the expression for consumption in equation (20), we obtain a formula for the model current account:


% Source PDF 23, text block 2.
\setcounter{equation}{22}
\begin{equation}\label{eq:23}
\begin{aligned}
CA_t=\frac{1}{1+\rho}\Biggl[&\underbrace{\left(\frac{D_t}{V_t}-\rho\right)\lambda_{t-1}V_t}_{\text{Equity in US}}+\underbrace{\left(\frac{D_t^*}{V_t^*}-\rho\right)\lambda_{t-1}^*V_t^*}_{\text{Equity in ROW}}\\
&+\underbrace{(r_t^*-\rho)B_t}_{\text{Bonds}}+\underbrace{\left(\frac{W_tL_t}{H_t}-\rho\right)H_t}_{\text{Human Wealth}}\Biggr].
\end{aligned}
\end{equation}


% Source PDF 23, text block 17.
This expression is intuitive. It compares the current income yield on each type of asset (both financial and human) owned by the household to the household’s rate of time preference and then takes a weighted aggregate of these quantities, where the weights are given by the beginning of period values of each type of asset held by the household. Domestic households will save out of dividend income on domestic and foreign equity if the current income yield on those assets exceeds their rate of time preference $\rho $. They will save out of bond income if the real interest rate exceeds $\rho $. And they will save out of labor income if the income yield on human wealth (current earnings relative to the future value of human wealth) exceeds $\rho $.


% Source PDF 23, text block 18.
This formula (23) reduces to a simple expression on a balanced growth path. Specifically, suppose the economy is on a balanced growth path, with a constant interest rate $\bar r^*$ and constant growth rate $\bar g$. On that path, the ratios $D_t/V_t$, $D_t^*/V_t^*$ and $W_tL_t/H_t$ are constant and equal to $(\bar r^*-\bar g)/(1+\bar g)$. It follows that if $\rho$ satisfies $1=\frac{1}{1+\rho}\frac{1+\bar r^*}{1+\bar g}$, then the balanced growth path current account will be $CA_t=\frac{1}{1+\rho}(\bar r^*-\rho)B_t=\bar g B_t$.

Off a balanced growth path, formula (23) captures the effects of business cycle shocks that lead to fluctuations in investment and corresponding fluctuations in $D_t$ and $D_t^*$. This formula also captures the effects of changes in the discount rate $r_{t+1}^*$ and expected future growth rate $\bar g_{t+1}$, which would also impact equilibrium current income yields on assets. When using our model for measurement, we match observed income yields on financial assets $D_t/V_t$ and $D_t^*/V_t^*$ from the data as described below. We derive how shocks impact the current account in our model in Section 4 of the Supplemental Appendix.


% Source PDF 23, text block 19.
The change in the end-of-period model net foreign asset position between $t-1$ and $t$ is given by the sum of the current account and asset revaluations:


% Source PDF 23, text block 20.
\setcounter{equation}{23}
\begin{equation}\label{eq:24}
\begin{aligned}
NFA_t-NFA_{t-1}=CA_t&+\lambda_{t-1}^*(V_t^*-V_{t-1}^*)\\
&-(1-\lambda_{t-1})(V_t-V_{t-1}),
\end{aligned}
\end{equation}


% Source PDF 23, text block 21.
where the final terms capture revaluations of foreign equity assets and liabilities. We refer to the term $\lambda _{t-1}^{*}(V_{t}^{*}-V_{t-1}^{*})$ as the revaluation of US equity assets in the ROW and the term $(1-\lambda _{t-1})(V_{t}-V_{t-1})$ as the revaluation of US equity liabilities to the ROW.


% Source PDF 23, text block 22.
The flow of capital must finance the current account, so we finally have


% Source PDF 23, text block 23.
\setcounter{equation}{24}
\begin{equation}\label{eq:25}
CA_t=B_{t+1}-B_t+(\lambda_t^*-\lambda_{t-1}^*)V_t^*-\left[(1-\lambda_t)-(1-\lambda_{t-1})\right]V_t.
\end{equation}



% ===== ORIGINAL PDF PAGE 24; JOURNAL PAGE 2174 =====
\sourcepage{24}{2174}

% Source PDF 24, text block 1.
\subsection*{G. Equilibrium}
\addcontentsline{toc}{subsection}{G. Equilibrium}


% Source PDF 24, text block 2.
An equilibrium is a set of sequences for the world interest rate $\{r_{t+1}^{*}\}_{t=0}^{\infty }$, for stock prices $\{V_{t},V_{t}^{*}\}_{t=0}^{\infty }$, for investment prices $\{Q_{t},Q_{t}^{*}\}_{t=0}^{\infty }$, and for domestic and foreign factor prices $\{R_{t},W_{t}\}_{t=0}^{\infty }$ and $\{R_{t}^{*},W_{t}^{*}\}_{t=0}^{\infty }$ such that when households and firms take these prices as given and solve their maximization problems with the expectations described above, all markets clear. Because bonds are in zero net supply, bond market clearing requires $B_{t+1}+B_{t+1}^{*}=0$.


% Source PDF 24, text block 3.
\section*{III. Using the Model for Measurement}
\addcontentsline{toc}{section}{III. Using the Model for Measurement}


% Source PDF 24, text block 4.
We use our model to measure the factors driving flows, stocks, and valuation of the US corporate sector and the US current account and net foreign asset position in two steps.


% Source PDF 24, text block 5.
In the first step, we use data from the Integrated Macroeconomic Accounts to construct model-consistent series for flows, stocks, and valuation of the US corporate sector and for the US current account and net foreign asset position. Specifically, we use the data in IMA Tables S5, S6, and S9 to construct quarterly series from 1990:I through 2023:III for corporate value added $Y_{t}$, taxes $\tau _{t}Y_{t}$, compensation of labor $W_{t}L_{t}$, consumption of fixed capital $\delta _{t}Q_{t}K_{t}$, investment expenditure $Q_{t}X_{t}$, end-of-period reproduction value of capital $Q_{t}K_{t+1}$, end-of-period enterprise value $V_{t}$, the current account $CA_{t}$, the NFA position $NFA_{t}$, US gross equity holdings in the ROW $\lambda _{t}^{*}V_{t}^{*}$, ROW gross holdings of equity in US corporations $(1-\lambda _{t})V_{t}$, as well as the gross flows and revaluations of these equity positions. We obtain data on monetary dividends paid on US equity in the ROW $\lambda _{t-1}^{*}D_{t}^{*}$ from the NIPA. We have reviewed much of these data in Section I. We provide detailed information on the data sources and construction of these variables in Appendix A.


% Source PDF 24, text block 6.
Note that these data imply measures of Gross and Net Operating Surplus and dividends from operations paid by US resident corporations as in equations (11), (13), and (14). Thus, these data summarize standard valuation metrics, including the ratio of the value of US resident corporations to GDP (the Buffett ratio) $V_{t}/Y_{t}$, the end of period reproduction value of capital to output ratio $Q_{t}K_{t+1}/Y_{t}$, Tobin’s Q measured as the ratio of the market valuation of the firm to the reproduction value of its capital stock $V_{t}/Q_{t}K_{t+1}$, the current dividend yields of US and ROW equity $D_{t}/V_{t},D_{t}^{*}/V_{t}^{*}$, and the US earnings yield, $E_{t}/V_{t}$.


% Source PDF 24, text block 7.
In the second step, we choose sequences of model parameters such that, as an equilibrium outcome, our model exactly reproduces the observed time series for all these data: the macro aggregates, the corporate valuations, and the current account. Specifically, we fix the rate of time preference for US households, $\rho $, to a constant value and solve analytically for sequences of parameters so that the model replicates the data items listed in the first step exactly for every quarter from 1990:I through 2023:III. The 12 time-varying parameters are (i) the discount rate for valuing the corporate sector $r_{t+1}^{*}$; (ii) the growth rate of aggregate productivity from $t+1$ on that is expected in period $t,\bar{g}_{t+1}$; (iii) the tax rate $\tau _{t}$; (iv) the depreciation rate $\delta _{t}$; (v–vi) domestic and foreign output wedges $\mu _{t}$ and $\mu _{t}^{*}$; (vii) labor’s share of costs $(1-\alpha _{t})$; (viii–ix) domestic and foreign replacement costs for capital $Q_{t}$ and $Q_{t}^{*}$; (x) the growth rate of productivity between $t$ and $t+1,g_{t+1}$; and (xi–xii) the gross cross-border equity positions $\lambda _{t}^{*}$ and $(1-\lambda _{t})$.



% ===== ORIGINAL PDF PAGE 25; JOURNAL PAGE 2175 =====
\sourcepage{25}{2175}

% Source PDF 25, text block 1.
We summarize the procedure for choosing these parameters here and provide a comprehensive explanation in Section  2 of the Supplemental Appendix.\hyperlink{fn:22}{\textsuperscript{22}} We describe how our measurement procedure relates to that used in prior macrofinance papers in detail in Section 3 of the Supplemental Appendix.


% Source PDF 25, text block 2.
We set the constant rate of time preference for US households $\rho $ to be consistent, on a balanced growth path, with the sample average of the current dividend yield on US corporations $D_{t}/V_{t}$.


% Source PDF 25, text block 3.
The evolution of the capital price $Q_{t}$, the depreciation rate $\delta _{t}$, the tax rate $\tau _{t}$, the productivity growth rate $g_{t+1}$ from $t$ to $t+1$, and foreign ownership of US equity $(1-\lambda _{t})$ are almost directly pinned down by data. In particular, the productivity growth rate $g_{t+1}$ is set equal to observed annualized quarterly growth in real US corporate value added, $\tau _{t}$ is set to match the share of corporate income going to taxes, and $\delta _{t}$ is set to match capital consumption. The sequence for $Q_{t}$ is chosen to ensure that equation (5), which links investment and depreciation flows to changes in the replacement value of capital, holds exactly at each date. The sequence for $(1-\lambda _{t})$ is simply the ratio of US equity liabilities relative to the enterprise value of the US corporate sector.


% Source PDF 25, text block 4.
Next, we turn to parameter values for the rest of the world. For our analysis, the only starred variables (besides $r_{t+1}^{*}$) that matter for the dynamics of US consumption (equation 20), the US current account (23), and the US net foreign asset position (24) are (i) US holdings of foreign equity, $\lambda _{t}^{*}$; (ii) foreign equity valuations, $V_{t}^{*}$; and (iii) foreign free cash flow, $D_{t}^{*}$. We set the path for $\lambda _{t}^{*}$ to replicate the path for US foreign equity assets relative to US corporate enterprise value. We set the series for $\mu _{t}^{*}$ and $Q_{t}^{*}$ to replicate the time series for $V_{t}^{*}$ and $D_{t}^{*}$. We assume that the time paths for all other rest-of-world parameters $(\alpha _{t}^{*},\delta _{t}^{*},\tau _{t}^{*},g_{t+1}^{*},\bar{g}_{t+1}^{*})$ are identical to their counterparts for the United States. Note that our choice to make $\mu _{t}^{*}$ and $Q_{t}^{*}$ the particular parameters whose time paths differ relative from their counterparts in the United States is arbitrary; any combinations of rest-of-world parameter sequences that replicate the observed paths for $V_{t}^{*}$ and $D_{t}^{*}$ will deliver identical time paths for all model observables.


% Source PDF 25, text block 5.
Four parameter series remain: $\alpha _{t},\mu _{t},r_{t+1}^{*}$, and $\bar{g}_{t+1}$. We set time paths for these four parameters to target four empirical times series: (i) labor’s share of income, $W_{t}L_{t}/Y_{t}$; (ii) end-of-period value of capital, $Q_{t}K_{t+1}$; (iii) the observed valuation of the corporate sector in excess of the replacement cost of capital, $V_{t}-Q_{t}K_{t+1}$; and (iv) the US current account. These four parameters are identified in a recursive and intuitive fashion.


% Source PDF 25, text block 6.
First, given a value for the rate of time preference $\rho $, all of the terms in the current account equation (23) are directly observed in data except for the value $H_{t}$ of human wealth from $t+1$ on, which is given in equation (22).\hyperlink{fn:23}{\textsuperscript{23}} This expression for human wealth is simply the ratio of labor compensation in period $t+1$ to $r_{t+1}^{*}-\bar{g}_{t+1}$. Because compensation is assumed to be known one period in advance, equations (22) and (23) can be used to directly identify $r_{t+1}^{*}-\bar{g}_{t+1}$. As a quantitative matter, human


% Source PDF 25, text block 7.
{\footnotesize \hypertarget{fn:22}{}\textbf{Footnote 22.} In particular, for simplicity, in the body of the paper, we have described the model in real terms. The data are in fact presented in nominal terms. We discuss the impact of inflation on our procedure for choosing parameters in Section 2 of the Supplemental Appendix.\par}


% Source PDF 25, text block 8.
{\footnotesize \hypertarget{fn:23}{}\textbf{Footnote 23.} This equation does involve the lagged value of the cost of capital $r_{t}^{*}$ and the stock of bonds carried into the period $B_{t}$. We construct series for $r_{t}^{*}$ and $B_{t}$ iteratively, using equations (23) and (25) to determine the stock of bonds carried into the next period, $B_{t+1}$. See Section 2 of the Supplemental Appendix for further details.\par}



% ===== ORIGINAL PDF PAGE 26; JOURNAL PAGE 2176 =====
\sourcepage{26}{2176}

% Source PDF 26, text block 1.
wealth $H_{t}$ is very large relative to the other components of wealth in equation (23). It follows that small fluctuations in $r_{t+1}^{*}-\bar{g}_{t+1}$ translate into large fluctuations in the current account. Because the actual current account is not particularly volatile, we will see that our quantification points to a near-constant time path for $r_{t+1}^{*}-\bar{g}_{t+1}$. Second, given $r_{t+1}^{*}-\bar{g}_{t+1}$, the time path for the output wedge parameter $\mu _{t+1}$ is identified from corporate valuations in excess of the replacement cost of capital, via equation (18). Recall that those valuations depend on $\mu _{t+1}$ via the expression for the expected share of profits in income, (8), which in turn generate asset valuations in proportion to the valuation multiple $1/(r_{t+1}^{*}-\bar{g}_{t+1})$.\hyperlink{fn:24}{\textsuperscript{24}}


% Source PDF 26, text block 2.
Third, given a path for $\mu _{t}$, the evolution of labor’s share of income identifies the path for $\alpha _{t}$ via equation (9).


% Source PDF 26, text block 3.
Finally, given values at each date for $\alpha _{t+1}$ and $\mu _{t+1}$, the evolution of the capital stock identifies the path for the discount factor $r_{t+1}^{*}$ via the firm’s first-order condition for investment, equation (12). In particular, given $\alpha _{t+1},\mu _{t+1},\tau _{t+1}$, and the observed capital output ratio at $t+1$, all of which are assumed to be known at $t$, agents can forecast the $t+1$ rental rate for capital, $R_{t+1}$. In addition, $\mathbb{E}_{t}[Q_{t+1}]=Q_{t}$, and $\delta _{t+1}$ is assumed known at $t$. Because all the components of the expected return to investment are known, the investment first-order condition can be used to identify the discount rate that rationalizes firms optimally choosing the observed end-of-period reproduction value of the capital stock $Q_{t}K_{t+1}$.


% Source PDF 26, text block 4.
There are parallels between this identification logic and the analysis of Lustig and Van Nieuwerburgh (2008). They start by noting that consumption growth is (i) much less volatile than stock returns and (ii) only weakly correlated with stock returns. They conclude that innovations to human wealth cannot be very volatile and must move inversely with innovations to financial wealth. In our model, output wedge shocks drive the lion’s share of valuation volatility. But those shocks naturally generate negative co-movement between financial and human wealth. And human wealth is not too volatile as long as fluctuations in $r^{*}-\bar{g}$ are small.

Note that in our model and measurement exercise, we abstract from value added created in the government sector and in the noncorporate private sector, which is important for residential housing. We have also abstracted from demographic factors relevant to the current account in an overlapping generations framework, as discussed in Auclert et al. (2021). To the extent that these omitted factors impact the US current account, they are implicitly captured in our model measurement in our estimates for $r_{t+1}^{*}-\bar{g}_{t+1}$ because the valuation multiple for human wealth is the only unknown in our equation (23) for the current account.


% Source PDF 26, text block 5.
To sum up, our parameterized model replicates exactly quarterly series for the following times series of the US corporate sector: value added, gross and net investment, labor earnings, taxes paid, cash flow payable to firm owners (defined as in equation (13)), the Buffett ratio, the replacement cost of capital, and the dividend and earnings yields. Figures 6 and 7 in Section 2 of the Supplemental Appendix plot these series.


% Source PDF 26, text block 6.
{\footnotesize \hypertarget{fn:24}{}\textbf{Footnote 24.} Note that, in our model, $\mu _{t+1}$ is determined by the ratio of the firm-specific productivity of the leader firm to that of the follower firm $z_{H,t+1}/z_{L,t+1}$ for each intermediate good. We set the firm-specific productivity $z_{H,t+1}$ each period so that US output in the model at $t+1$ is equal to $Z_{t+1}=(1+g_{t+1})Z_{t}$. We will follow the same procedure for the ROW. Thus, the values for model output in the US and the ROW are identical at each date $t$.\par}



% ===== ORIGINAL PDF PAGE 27; JOURNAL PAGE 2177 =====
\sourcepage{27}{2177}

% Figure 10: original PDF page 27, plot crop x=60..444, y=62..371 points.
\begin{center}
\begingroup\renewcommand{\thefigure}{10}\refstepcounter{figure}\label{fig:10}\endgroup
\includegraphics[page=27,trim=60 349 60 62,clip,width=0.95\linewidth]{\sourcepdf}
\par\small Figure 10. Estimated Time Series for Key Parameters
\end{center}


% Source PDF 27, text block 1.
\section*{IV. Baseline Results}
\addcontentsline{toc}{section}{IV. Baseline Results}


% Source PDF 27, text block 2.
Figure 10 plots some of the key parameter sequences derived from our measurement procedure.\hyperlink{fn:25}{\textsuperscript{25}} The top left panel plots (in blue) the model-implied sequence for $r_{t}^{*}$ (annualized), while the top right panel shows the sequence for trend growth $\bar{g}_{t}$ alongside actual quarterly growth $g_{t}$, both also annualized. The estimated sequence for trend growth is much less volatile than the actual quarterly growth rate series, but some correlation between the two is apparent. The sequence for $r_{t}^{*}$ shown in Figure 10 declines substantially over the past decade. However, our measurement procedure implies that the sequence for the discount rate $r_{t}^{*}$ tracks the trend growth rate series closely, so that the difference between the two (the red line in the top left panel of Figure 10) fluctuates very little.


% Source PDF 27, text block 3.
The bottom left panel shows the model-implied time series for the output wedge for the US corporate sector, $\mu _{t}$; the path for share of factorless income $(1-\tau _{t})(\mu _{t}-1)/\mu _{t}$ looks similar. The share of capital in costs, $\alpha _{t}$, shown in the bottom right panel, exhibits some fluctuations that are inversely correlated with those in $\mu _{t}$ but features no clear long-run trend.


% Source PDF 27, text block 4.
{\footnotesize \hypertarget{fn:25}{}\textbf{Footnote 25.} Figure 5 in Section 2 of the Supplemental Appendix plots time series for all parameter values.\par}



% ===== ORIGINAL PDF PAGE 28; JOURNAL PAGE 2178 =====
\sourcepage{28}{2178}

% Source PDF 28, text block 1.
To understand why the model measurement procedure produces these time paths for parameters, recall the logic of the recursive identification scheme outlined above.


% Source PDF 28, text block 2.
First, the reason that we find a stable valuation multiple for factorless income is rooted in our requirement that the model match the observed current account sequence for the United States. The current account is highly sensitive to changes in the ratio of current labor income to human wealth because the value of human wealth, $H_{t}$, is large. Moreover, $H_{t}$ is inversely proportional to $r_{t+1}^{*}-\bar{g}_{t+1}$ (see equation (22)). Thus, large changes in $r_{t+1}^{*}-\bar{g}_{t+1}$ would imply large changes in desired US consumption and counterfactually large swings in the current account. Put differently, the observed current account is a data moment that provides sharp identification for expected trend growth.\hyperlink{fn:26}{\textsuperscript{26}}


% Source PDF 28, text block 3.
Second, consider first the model’s implications for changes in the discounted present value of firms’ future factorless income, $V_{t}^{\Pi }$. Because the valuation multiple $(r_{t+1}^{*}-\bar{g}_{t+1})^{-1}$ applied to expected future factorless income is quite stable, the model attributes the majority of the observed fluctuations in this component of firms’ valuation to fluctuations in the quantity of factorless income $\Pi _{t+1}$. Changes in measured taxes explain some of the model-inferred path for factorless income, but the rest is attributed to latent changes in the output wedge, $\mu _{t}$. Thus, the identified series for the output wedge $\mu _{t}$ closely tracks the $V_{t}^{\Pi }$ component of valuation. For example, the output wedge spikes in the tech boom of 2000, when valuations were temporarily elevated. Note that because the valuation multiple for factorless income is quite stable, the decline in the discount rate that the model identifies in the post-2010 period does not account for much of the observed run-up in asset valuations in excess of the replacement cost of capital over that period: Future expected profits are discounted at a lower rate but are also expected to grow more slowly, with the two effects offsetting.


% Source PDF 28, text block 4.
Our series for the output wedge $\mu _{t}$ can be compared to other estimates for profitability, such as Figure 8 in De Loecker, Eeckhout, and Unger (2020), Figure 3 in Barkai (2020), or Figure 7 in Atkeson, Heathcote, and Perri (2025). All show an increase, of around 10 percentage points, in the share of factorless income since 1990, albeit with somewhat different timing.\hyperlink{fn:27}{\textsuperscript{27}} We view the fact that all estimates of the overall size of the increase are of similar magnitude as strengthening our finding that the share of factorless income in the United States has indeed increased; we also believe that understanding the differences in the timing of the increase is an interesting direction for further research.


% Source PDF 28, text block 5.
Third, given the path for $\mu _{t}$ that replicates the observed path for the value of claims to factorless income, $\alpha _{t}$ must fluctuate to match the observed path for labor’s share of income. If labor’s share is not moving, $\alpha _{t}$ must move inversely with $\mu _{t}$ so that a rise in profit’s share of income is offset by a decline in capital’s share. That


% Source PDF 28, text block 6.
{\footnotesize \hypertarget{fn:26}{}\textbf{Footnote 26.} This logic is also present in Aguiar and Gopinath (2007), who argue that the current account swings in emerging markets can be used to identify changes in expected growth.\par}


% Source PDF 28, text block 7.
{\footnotesize \hypertarget{fn:27}{}\textbf{Footnote 27.} De Loecker, Eeckhout, and Unger (2020) estimate profit share income using microdata on markups. Barkai (2020) estimates profit’s share of value added by imputing a series for the rate of return on capital given which he can infer capital costs. Profit’s share under his approach is value added minus the sum of compensation of employees, capital costs, and taxes. Finally Atkeson, Heathcote, and Perri (2025) estimate profit share assuming a constant division of income between capital and labor so that increase in profit share is directly identified from the decline in labor share.\par}



% ===== ORIGINAL PDF PAGE 29; JOURNAL PAGE 2179 =====
\sourcepage{29}{2179}

% Source PDF 29, text block 1.
inverse co-movement is apparent in Figure 10. For example, as $\mu _{t}$ rises around 2000 to replicate the tech stock boom, $\alpha _{t}$ simultaneously falls so that labor’s share is relatively stable, as in the data. In the 2000s, $\alpha _{t}$ rises to allow the model to replicate a declining labor share. After 2010, labor’s share is relatively flat, but the estimated profit share is rising, so $\alpha _{t}$ must decline to offset a rising path for $\mu _{t}$.


% Source PDF 29, text block 2.
Finally, consider the path for $r_{t}^{*}$. Recall that this path is such that the model reproduces the observed path for the replacement value of installed capital in US resident corporations. That value has been quite stable over time, relative to corporate value added (see the left panel of Figure 8). That stability emerges naturally in the prototypical stochastic growth model in which all structural parameters are constant, and fluctuations are driven by transitory productivity shocks. But in our model, shocks to $\alpha _{t+1}$ and $\mu _{t+1}$ translate into fluctuations in the expected rental rate for capital. All else equal, those shocks would translate into changes in firms’ desired capital stock, relative to value added. Thus, the model infers offsetting changes in the discount rate. For example, in the post-2010 period, a declining path for $\alpha _{t}$ and a rising output wedge $\mu _{t}$ both work to depress expected rental income from capital. But the actual replacement value of capital relative to value added did not fall over this period. The model rationalizes that finding by inferring that the discount rate used by investors must have declined.


% Source PDF 29, text block 3.
Note that this logic reframes the conventional narrative on the relationship between investment and interest rates. That narrative has been that investment has been surprisingly weak in recent decades given low estimates for the cost of capital (see, e.g., Gutiérrez and Philippon 2023). In the context of our model, in contrast, investment appears surprisingly strong, given a declining share of income accruing to capital. The model makes sense of this surprisingly strong investment by inferring a declining cost of capital.


% Source PDF 29, text block 4.
One interesting observation from valuation data is that while the US corporate dividend yield exhibits no long-term trend, there is a downward trend in the earnings yield (see Figure 7 in Section 2 of the Supplemental Appendix). What explains this difference between these two standard valuation metrics? In our model, the expected forward dividend yield is $(r_{t+1}^{*}-\bar{g}_{t+1})$, as in the Gordon growth formula:


% Source PDF 29, text block 5.
\setcounter{equation}{25}
\begin{equation}\label{eq:26}
r_{t+1}^*-\bar g_{t+1}=\frac{\mathbb{E}_t[D_{t+1}]}{V_t}.
\end{equation}


% Source PDF 29, text block 6.
We define earnings similarly to dividends, except that depreciation is subtracted from gross operating surplus instead of investment. Thus, the expected forward earnings to price ratio is also $(r_{t+1}^{*}-\bar{g}_{t+1})$ for monopolists, for whom there is no distinction between earnings and dividends. But for investment firms, the expected earnings yield is $r_{t+1}^{*}$. The expected earnings yield for the entire economy is a weighted average of the ratios for the two firm types, with weights given by their shares in total enterprise value:


% Source PDF 29, text block 7.
\[
\frac{\mathbb{E}_t[E_{t+1}]}{V_t}=\left(r_{t+1}^*-\bar g_{t+1}\right)\frac{V_t^{\Pi}}{V_t}+r_{t+1}^*\frac{Q_tK_{t+1}}{V_t}.
\]


% Source PDF 29, text block 8.
Note that when trend growth is positive, the forward earnings yield will exceed the forward dividend yield (as is evident in Figure C.3). The intuition is that in a growing economy, investment exceeds depreciation, so cash flow payable to investors is



% ===== ORIGINAL PDF PAGE 30; JOURNAL PAGE 2180 =====
\sourcepage{30}{2180}

% Figure 11: original PDF page 30, plot crop x=60..444, y=62..364 points.
\begin{center}
\begingroup\renewcommand{\thefigure}{11}\refstepcounter{figure}\label{fig:11}\endgroup
\includegraphics[page=30,trim=60 356 60 62,clip,width=0.95\linewidth]{\sourcepdf}
\par\small Figure 11. Current Account Decomposition Following Equation (23)
\end{center}


% Source PDF 30, text block 1.
less than earnings. This equation also explains why the expected earnings yield in model and data declines over time, even absent trends in $r_{t+1}^{*}$ or $\bar{g}_{t+1}$. In particular, as the share of monopolist firms in total enterprise value rises over time, the aggregate expected earnings yield puts more weight on the lower value $(r_{t+1}^{*}-\bar{g}_{t+1})$ for monopolist firms and less weight on the higher value $r_{t+1}^{*}$ for investment firms.

Figure 11 plots one of the main results in our paper. The top left panel is the US current account relative to corporate sector value added, which the model is calibrated to perfectly replicate. The US current account deficit widened steadily through the 1990s and early 2000s, before moderating during the Great Recession. The other panels of the figure decompose the model current account series following equation (23). The decline in the US current account during the 1990s is primarily attributed to declining dividend yields on US equity, which led US households to borrow from the rest of the world. However, when US asset values fall during the dot-com bust, that dividend yield rises, which, all else equal, would have pushed the current account back toward balance. The model rationalizes the continuing observed decline in the current account via rising expected growth (see Figure 10), which depresses $r_{t+1}^{*}-\bar{g}_{t+1}$, raising the value of human wealth and stimulating ongoing borrowing. A high dividend yield on US equity during the Great Recession helps explain why the current account narrowed around this time. Thereafter, slowing expected growth reduces the value of human capital, which is why current account deficits remain modest even as the dividend yield on US equity declines.



% ===== ORIGINAL PDF PAGE 31; JOURNAL PAGE 2181 =====
\sourcepage{31}{2181}

% Figure 12: original PDF page 31, plot crop x=60..444, y=62..372 points.
\begin{center}
\begingroup\renewcommand{\thefigure}{12}\refstepcounter{figure}\label{fig:12}\endgroup
\includegraphics[page=31,trim=60 348 60 62,clip,width=0.95\linewidth]{\sourcepdf}
\par\small Figure 12. NFA Dynamics
\end{center}


% Source PDF 31, text block 1.
The top left panel of Figure 12 plots the NFA position predicted by the model against the actual position. The top right panel plots the net equity position, which matches the data by construction. The NFA position in the model reflects cumulative current accounts (bottom left panel) plus cumulative equity revaluations (bottom right panel).


% Source PDF 31, text block 2.
There are three reasons why the model does not perfectly replicate the data path for the NFA position. The first is that the equity liability revaluations in the model are not identical to those reported in Table S9 in the Integrated Macroeconomic Accounts. Recall that we infer our own series for the revaluation of ROW holdings of US equity based on our estimated series for the value of the US corporate sector. Note, however, that the difference between the two revaluation series is small (bottom right panel of Figure 12). The second reason is that the model does not feature any non-equity valuation effects, while there are such effects in the data (see Figure 4). The third and most important reason is that in the data accounting identity (equation (1)), there is a residual term that is not in the model. This residual term, which incorporates the statistical discrepancy between the current account and net foreign asset purchases, contributed to a significant improvement in the US NFA position in the 2000s, which our model cannot replicate (see Figure 2).\hyperlink{fn:28}{\textsuperscript{28}}


% Source PDF 31, text block 3.
{\footnotesize \hypertarget{fn:28}{}\textbf{Footnote 28.} The combined impact of non-equity revaluations and the residual term can be seen by comparing the cumulative actual current account in the data (which the model perfectly replicates) to the hypothetical current account\par}



% ===== ORIGINAL PDF PAGE 32; JOURNAL PAGE 2182 =====
\sourcepage{32}{2182}

% Source PDF 32, text block 1.
Expected versus Unexpected Shocks.—Our focus is on the impact of changes in asset values and returns on the current account and the NFA position. Some portion of the changes in asset valuations seen in the data are due simply to anticipated factors; as an economy grows and as it invests in more physical capital, the value of its corporations would be expected to grow as well. Other changes in asset valuations are due to the arrival of news that makes realized returns differ from expected returns. We now decompose changes in the US NFA position into anticipated and unanticipated factors.


% Source PDF 32, text block 2.
Given the information available to households in period $t-1$ about model parameters, they expect returns on all assets between $t-1$ and $t$ to be equal to $r_{t}^{*}$. Then, at each date $t$, these households receive further news about future output wedges, growth, and other parameters, and this news generates dynamics both in the current account and in asset values. Specifically, let $e_{t}$ and $e_{t}^{*}$ denote excess real returns (realized minus expected return) to domestic and foreign equity in period $t$:


% Source PDF 32, text block 3.
\[
e_t=\frac{D_t+V_t}{V_{t-1}}-(1+r_t^*),\qquad e_t^*=\frac{D_t^*+V_t^*}{V_{t-1}^*}-(1+r_t^*).
\]


% Source PDF 32, text block 4.
In the model, these excess returns are due to news about parameters from $t+1$ onward that leads agents to expect a different present value of dividend income at $t$ relative to what they expected at $t-1$.


% Source PDF 32, text block 5.
In Section 6 of the Supplemental Appendix, we show that the evolution of the US net foreign asset position in our model can be decomposed as follows:


% Source PDF 32, text block 6.
\setcounter{equation}{26}
\begin{equation}\label{eq:27}
\begin{aligned}
NFA_t-NFA_{t-1}={}&\underbrace{\frac{r_t^*-\rho}{1+\rho}NFA_{t-1}}_{(1)}
+\underbrace{\left(\frac{r_t^*-\rho}{1+\rho}-\bar g_t\right)V_{t-1}}_{(2)}\\
&+\underbrace{\left(\frac{W_tL_t/H_t-\rho}{1+\rho}\right)H_t}_{(3)}
-\underbrace{\left(Q_tX_t-\mathbb{E}_{t-1}[Q_tX_t]\right)}_{(4)}\\
&-\underbrace{\frac{\rho}{1+\rho}\lambda_{t-1}e_tV_{t-1}}_{(5)}
-\underbrace{\frac{\rho}{1+\rho}\lambda_{t-1}^*e_t^*V_{t-1}^*}_{(6)}\\
&-\underbrace{e_t(1-\lambda_{t-1})V_{t-1}}_{(7)}
+\underbrace{e_t^*\lambda_{t-1}^*V_{t-1}^*}_{(8)}.
\end{aligned}
\end{equation}


% Source PDF 32, text block 31.
Terms (1) and (2) capture savings motives that are predictable at $t-1$. Term (1) indicates that the NFA position will tend to grow at rate $(r_{t}^{*}-\rho )/(1+\rho )$. Note that on a balanced growth path, this term is equal to the growth rate $\bar{g}$, implying a stable NFA to GDP ratio.\hyperlink{fn:29}{\textsuperscript{29}} The second term is the contribution to the NFA from expected returns to domestic equity. Term (3) is the current account contribution to national saving from a yield on human capital exceeding $\rho $.


% Source PDF 32, text block 32.
{\footnotesize\textbf{Footnote 28 (continued).} series that would obtain in the data absent a residual term and absent non-equity revaluations. That hypothetical series is plotted in yellow in the bottom left panel of Figure 12.\par}


% Source PDF 32, text block 33.
{\footnotesize\hypertarget{fn:29}{}\textbf{Footnote 29.} In particular, if $1=\frac{1}{1+\rho}\frac{1+r^*}{1+\bar g}$, then $\frac{r^*-\rho}{1+\rho}=\bar g$.\par}



% ===== ORIGINAL PDF PAGE 33; JOURNAL PAGE 2183 =====
\sourcepage{33}{2183}

% Figure 13: original PDF page 33, plot crop x=60..444, y=62..295 points.
\begin{center}
\begingroup\renewcommand{\thefigure}{13}\refstepcounter{figure}\label{fig:13}\endgroup
\includegraphics[page=33,trim=60 425 60 62,clip,width=0.95\linewidth]{\sourcepdf}
\par\small Figure 13. NFA Decomposition
\end{center}


% Source PDF 33, text block 1.
The remaining terms capture the impact of shocks at $t$ to asset values and returns. Term (4) captures deviations of investment at date $t$ from investment expected at $t-1$: If information revealed at $t$ spurs unexpected domestic investment, the United States will fund that extra investment by borrowing from abroad. Terms (5) and (6) capture the impact of excess equity returns at $t$ on desired consumption and thus the current account. Terms (7) and (8) capture the direct effect of excess returns on the NFA position: Here, excess returns to domestic equity reduce the NFA position by inflating US liabilities, while excess returns to foreign equity improve the position. On a balanced growth path, terms (2) through (8) are all zero.


% Source PDF 33, text block 2.
Figure 13 uses equation (27) to decompose the change in the NFA position relative to the start of our sample period (1990) into the cumulative values of each of the eight labeled terms. The cumulative value for each term is plotted relative to value added at date $t$. The message from the plot is that over the entire sample period, there are three key drivers of the decline in the US NFA position.


% Source PDF 33, text block 3.
Quantitatively, the most important driver is that starting around the end of the Great Recession, positive excess returns on US equities have inflated US equity liabilities (shown in the purple line in the right panel of Figure 13). These excess returns account for a decline in the NFA position of over 100 percent of corporate value added, though some of this effect was unwound in 2022 as US asset values declined. Note that, at high frequency, excess returns to domestic and foreign equity tend to work in opposite directions, reflecting high-frequency co-movement across global equity markets. For example, in 2022 we see that poor US equity market performance reduced US equity liabilities (purple line), but poor ROW market performance simultaneously reduced the value of US foreign equity assets (green line).



% ===== ORIGINAL PDF PAGE 34; JOURNAL PAGE 2184 =====
\sourcepage{34}{2184}

% Source PDF 34, text block 1.
The second important driver of the US NFA position is that during the 2000s, a low income yield on human capital fueled current account deficits (shown in the yellow line in the left panel of Figure 13.) The third key driver is that as the NFA position has widened, the fact that the interest rate exceeds the rate of time preference has fueled further borrowing (the blue line in the left panel).


% Source PDF 34, text block 2.
\section*{V. Counterfactuals}
\addcontentsline{toc}{section}{V. Counterfactuals}


% Source PDF 34, text block 3.
In Section III, we used our model to measure the factors driving observed flows, stocks, and valuations of the US corporate sector, together with those driving the evolution of the US current account and NFA position in quarterly data over the period 1990 through 2023. This measurement exercise establishes a baseline path for model parameters that accounts for the evolution of this broad collection of data over this three-decade time period. We now use the model to conduct counterfactual exercises relative to this model baseline to consider how these changes in parameters impacted the welfare of US households. Our particular focus is on the question of how the large increase in gross cross-border equity positions observed in recent decades impacts the welfare implications of the large increase in the output wedge $\mu _{t}$ measured in our baseline over the course of the last decade. To address this question, we simulate the model for three specific counterfactual paths for parameters.


% Source PDF 34, text block 4.
In the first, we solve for the counterfactual equilibrium of the model with no cross-border equity positions. To do so, we set the share of US firms owned by US residents to $\lambda _{t}=1$ and the share of ROW firms owned by US residents to $\lambda _{t}^{*}=0$, leaving all other parameter values unchanged. In this counterfactual exercise, the paths of output, labor compensation, capital, investment, and the market valuation of the US corporate sector are all the same as in the baseline calibration. That is because, in our model, the distribution of the ownership of firms across US and ROW households does not impact the equilibrium discount rate $r_{t}^{*}$ or equilibrium production and investment decisions. In this counterfactual exercise, the paths for US consumption and for the current account are different from those in the baseline because now US households enjoy more of the unexpected capital gains on their now-larger holdings of equity in US corporations. Our baseline model, together with this counterfactual exercise, gives us two paths for the consumption of US households with a large increase in $\mu _{t}$: one with and one without observed gross cross-border equity positions.


% Source PDF 34, text block 5.
In the second counterfactual exercise, we solve for the equilibrium of the model when $z_{Lt}$ is set equal to $z_{Ht}$ for all $t$, leaving all other parameter values unchanged. In this counterfactual exercise, with this alternative path for the productivity of follower firms, we have $\mu _{t}=1$ for all $t$. This counterfactual exercise gives us predictions for how allocations would have evolved if US leader firms had experienced the same path for productivity but had not enjoyed a large increase in their power to price above cost, all under the assumption that cross-border equity share holdings $(1-\lambda _{t})$ and $\lambda _{t}^{*}$ had evolved as in the baseline.

In the third counterfactual exercise, we solve for the equilibrium path for consumption of US households when we set the share of US firms owned by US residents to $\lambda _{t}=1$, the share of ROW firms owned by US residents to $\lambda _{t}^{*}=0$, and $z_{Lt}=z_{Ht}$ for all $t$ so that $\mu _{t}=1$ for all $t$.



% ===== ORIGINAL PDF PAGE 35; JOURNAL PAGE 2185 =====
\sourcepage{35}{2185}

% Figure 14: original PDF page 35, plot crop x=60..444, y=62..374 points.
\begin{center}
\begingroup\renewcommand{\thefigure}{14}\refstepcounter{figure}\label{fig:14}\endgroup
\includegraphics[page=35,trim=60 346 60 62,clip,width=0.95\linewidth]{\sourcepdf}
\par\small Figure 14. Effect of Output Wedges on $Y,K/Y$, and $C$\par Note: Effect on $C$ shown with actual path for diversification and zero diversification counterfactual.
\end{center}


% Source PDF 35, text block 1.
With these three counterfactuals, we can compare the paths of output and consumption given the baseline path for the output wedge $\mu _{t}$ against the paths of output and consumption with $\mu _{t}=1$. And we can perform this comparison twice: once under the baseline data-consistent paths for gross equity positions $(1-\lambda _{t})$ and $\lambda _{t}^{*}$, and once for a counterfactual world in which cross-border equity positions are always zero. We use these counterfactual simulations to study how international equity diversification changes the welfare implications of a rise in the ratio between revenue and cost $\mu _{t}$, which is the key driver of rising US asset valuations in our analysis.


% Source PDF 35, text block 2.
We show these results in Figure 14. In the top left panel of this figure, we show the ratio of the path of US output in our baseline to that in our counterfactuals with $z_{Lt}=z_{Ht}$ for all $t$ so that $\mu _{t}=1$ for all $t$. As noted above, the path for output implied by our model is not impacted by assumptions about the extent of international diversification of equity holdings $(1-\lambda _{t})$ and $\lambda _{t}^{*}$. Hence, there is only one line in this top panel. We see in this figure that the increase in $\mu _{t}$ over the past decade has had a large negative impact on the path of output relative to the counterfactual with $\mu _{t}=1$ at all dates. In our model, since labor is supplied inelastically and the path for leader firm productivity is the same in the baseline and the counterfactual



% ===== ORIGINAL PDF PAGE 36; JOURNAL PAGE 2186 =====
\sourcepage{36}{2186}

% Source PDF 36, text block 1.
simulation, this decline in output is entirely due to the impact of increases in $\mu _{t}$ on the accumulation of physical capital; the negative impact on the model capital to output ratio is plotted in the top right panel of the figure. In this regard, we confirm the findings of Gutiérrez and Philippon (2017) that the big increase in the valuation of US firms in our baseline is associated with comparatively weak investment.


% Source PDF 36, text block 2.
In the bottom panels of Figure 14, we show how the impact of a rising output wedge on the consumption (and thus welfare) of US households is mediated by the dynamics of international equity diversification.


% Source PDF 36, text block 3.
First, consider the impact of the large estimated increase in the output wedge $\mu _{t}$ in a world with no cross-border equity holdings (bottom right panel). Here, we plot the ratio of the path of US household consumption in our baseline (in which $\mu _{t}$ is generally rising) to consumption in the counterfactual with $z_{Lt}=z_{Ht}$ for all $t$ (so that $\mu _{t}=1$ for all $t$) under the assumption that cross-border equity holdings $(1-\lambda _{t})$ and $\lambda _{t}^{*}$ are always zero in both model simulations. It is clear that a large increase in the output wedge $\mu _{t}$ in the United States has only a modest impact on US consumption, notwithstanding the large impact of this increase in the output wedge on US output shown in the top-left panel. The intuition for this result is straightforward. First, the income that US households lose from lower labor compensation when $\mu _{t}$ goes up is mostly offset by a rise in the factorless income that they receive in dividend payments from firms. Second, the decline in investment that drives down equilibrium output also implies a period of elevated cash flow to shareholders, and that extra income can be invested abroad and used to replace lower future domestic income. The finding that a rise in $\mu _{t}$ has only a small negative impact on consumption is analogous to the result that an increase in output wedges starting from an efficient equilibrium has no first-order impact on welfare in a closed economy. In Appendix B we prove that if we start from a zero output wedge, the impact on consumption from a marginal shock to $\mu _{t}$ is nil when there is zero foreign ownership of US equity.


% Source PDF 36, text block 4.
Now consider the impact of the same increase in $\mu _{t}$ under our baseline parameterization, in which cross-border equity holdings match those in the data. Specifically, the bottom left panel of Figure 14 plots the path for US household consumption in our baseline relative to that in our counterfactual with $z_{Lt}=z_{Ht}$ for all $t$ (so that $\mu _{t}=1$) under the assumption that cross-border equity holdings $(1-\lambda _{t})$ and $\lambda _{t}^{*}$ evolve as in our baseline calibration. Now, we see a large negative impact of the increase in $\mu _{t}$ on equilibrium consumption. In fact, the additional decline in consumption relative to a world with a zero output wedge is quantitatively similar to the associated decline in output. The reason is that now the income that US households lose from a decline in the share of labor compensation in output is not largely offset by an increase in dividends that they receive as owners of firms because foreign households receive a large share of those increased dividends. As a result, US households have less wealth (both in absolute terms and relative to output) than they would have had absent cross-border equity holdings, leading them to reduce consumption. Thus, the negative welfare implications of an increase in US firms’ market power for US households are dramatically magnified in the presence of large foreign ownership of US equity. In the alternative calibration to smaller estimates for gross cross-border equity positions described in Section 1 of the Supplemental Appendix, we find slightly smaller welfare losses for US households, with a maximum consumption decline relative to the $\mu _{t}=1$ counterfactual of 4 percent rather than 5 percent.



% ===== ORIGINAL PDF PAGE 37; JOURNAL PAGE 2187 =====
\sourcepage{37}{2187}

% Source PDF 37, text block 1.
Diversification.—Is the ex post redistribution to foreign households that occurs when US factorless income increases following the Global Financial Crisis desirable from an ex ante risk-sharing perspective?


% Source PDF 37, text block 2.
In the model we have explored to this point, agents perceive zero risk. Thus, as long as domestic equity, foreign equity, and bonds offer the same expected return, agents are indifferent about their portfolio choices. This indifference allowed us to focus on an equilibrium in which portfolios match those observed in the data.


% Source PDF 37, text block 3.
We now consider a model extension in which households perceive risk and in which risk-sharing considerations dictate a unique portfolio choice. We consider two alternative cases: one in which there are fluctuations in $\mu _{t+1}$ that do not impact aggregate income $Y_{t+1}$, and a second in which there are conventional shocks to productivity $Z_{t+1}$ that translate into fluctuations in $Y_{t+1}$ without impacting the output wedge $\mu _{t+1}$.\hyperlink{fn:30}{\textsuperscript{30}} We abstract here from all other shocks and retain the assumption that foreign investors are risk neutral. Thus, bonds and foreign equity pay a constant return $r^{*}$, and agents expect constant future productivity growth at rate $\bar{g}$. We will assume $1+\rho =(1+r^{*})/(1+\bar{g})$.

Consider first the experiment with shocks to $\mu _{t+1}$. We model risk as follows. Let $\kappa _{t}=(\mu _{t}-1)/\mu _{t}$ denote the share of after-tax income accruing to monopoly profits. At each date $t$ agents receive news about $\kappa _{t+1}$:


% Source PDF 37, text block 4.
\[
\kappa_{t+1}=\kappa_t+\varepsilon_{\kappa,t},
\]


% Source PDF 37, text block 5.
where the news shock $\varepsilon _{\kappa ,t}$ is drawn from a distribution with mean zero. These news shocks have two implications for domestic agents. First, a positive value for $\varepsilon _{\kappa ,t}$ signals higher future profits and will result in an immediate increase in the value of domestic equity. Second, the same shock signals a decline in labor’s share of after-tax income, $(1-\kappa _{t})(1-\alpha )$, and thus lower future labor income. In this environment, what portfolios are chosen in equilibrium?


% Source PDF 37, text block 6.
\hypertarget{prop:1}{}\textbf{PROPOSITION 1:} Consider the environment just described. At date 0, for any initial state, equilibrium allocations have the following properties: (i) $\lambda _{t}=(1-\alpha )$, and (ii) $C_{t+1}/C_{t}=1+\bar{g}$ for all $t\geq 0$.


% Source PDF 37, text block 7.
\textbf{PROOF:}


% Source PDF 37, text block 8.
See Appendix B.


% Source PDF 37, text block 9.
The logic for this result is that a news shock that increases the expected share of income accruing to owners of firms by 1 percentage point reduces the share of income accruing to labor by $(1-\alpha )$ percentage points. When domestic households own exactly fraction $(1-\alpha )$ of domestic firms, the shock has no impact on their expected flow income or on the sum of their human plus financial wealth. Thus, this portfolio delivers perfect insurance, in that domestic consumption grows at the trend growth rate and is independent of the realized path for


% Source PDF 37, text block 10.
{\footnotesize \hypertarget{fn:30}{}\textbf{Footnote 30.} In the first case, we assume that changes in $\mu _{t+1}$ reflect changes in the ratio of productivities of leader to follower firms $z_{H,t+1}/z_{L,t+1}$, with the property that equilibrium output at $t+1$ is unaffected; see footnote 24 and equation (6) in Section 2 of the Supplemental Appendix.\par}



% ===== ORIGINAL PDF PAGE 38; JOURNAL PAGE 2188 =====
\sourcepage{38}{2188}

% Source PDF 38, text block 1.
$\kappa _{t}$. Equilibrium portfolios are home-biased (assuming $\alpha <0.5$) because shocks that reduce labor income will simultaneously generate excess returns to domestic equity. This mirrors earlier results in the literature that home-biased equity portfolios provide insurance against redistributive shocks; see, e.g., Coeurdacier, Kollmann, and Martin (2007).


% Source PDF 38, text block 2.
Compare this result to equilibrium portfolio choice in a conventional model in which $\kappa $ is constant but in which there are news shocks about output growth. In particular, suppose that


% Source PDF 38, text block 3.
\[
Z_{t+1}=(1+\bar g)\varepsilon_{Z,t}Z_t,
\]


% Source PDF 38, text block 4.
where the news shock $\varepsilon _{Z,t}$ is drawn from a distribution with mean one.


% Source PDF 38, text block 5.
\hypertarget{prop:2}{}\textbf{PROPOSITION 2:} Consider the environment just described. At date 0, for any initial state, equilibrium allocations have the following properties: (i) $\lambda _{t}=-(1-\alpha )(1-\kappa )/\kappa $, and (ii) $C_{t+1}/C_{t}=1+\bar{g}$ for all $t\geq 0$.


% Source PDF 38, text block 6.
\textbf{PROOF:}


% Source PDF 38, text block 7.
See Appendix B.


% Source PDF 38, text block 8.
Thus, if risk takes the form of conventional uncertainty about the path for country-specific TFP, the model predicts that households should short domestic equity. The logic is that, in this case, shocks to domestic labor income are positively correlated with returns to domestic equity, so a large short position is required to insulate households against those shocks (see also Baxter and Jermann 1997).


% Source PDF 38, text block 9.
To summarize, in economies with risk, country-specific shocks generate country-specific excess returns to equity. The pattern of international equity diversification determines how those excess returns impact the net foreign asset position. In equilibrium, equity portfolio choice is designed to provide insurance against labor income risk, which cannot be insured directly.


% Source PDF 38, text block 10.
We conclude this section by returning to the question we started with. In the period following the Great Recession, our calibration attributes excess returns to US equity to unexpected increases in the share of factorless income accruing to the owners of US firms, which translated into ex post redistribution to foreign investors. Can those transfers be interpreted as part of an ex ante optimal risk sharing arrangement? Interestingly, our estimated value over this period for the labor exponent in production, $1-\alpha _{t}$, is close to our measure of the share of US equity that is domestically owned, $\lambda _{t}$ (see Figure 5 in Section 2 of the Supplemental Appendix). Thus, if the key macro risk investors worried about over this period was risk to the factorless income share $\kappa _{t}$, then Proposition 1 suggests that diversification has perhaps been close to optimal, and ex post transfers to foreign shareholders may have been ex ante efficient. However, a complete quantitative assessment of the optimal extent of diversification for US households is beyond the scope of this paper. Such an assessment would require modeling the complete set of shocks that drive aggregate fluctuations and specifying the joint stochastic process for those shocks, including how domestic shocks co-move with their foreign counterparts (see, for example, Coeurdacier and Rey 2013 and Heathcote and Perri 2013).



% ===== ORIGINAL PDF PAGE 39; JOURNAL PAGE 2189 =====
\sourcepage{39}{2189}

% Source PDF 39, text block 1.
\section*{VI. Sensitivity Analysis}
\addcontentsline{toc}{section}{VI. Sensitivity Analysis}


% Source PDF 39, text block 2.
We now consider the extent to which our measurement of the discount rate $r_{t+1}^{*}$ and the parameters $\mu _{t+1}$ and $\alpha _{t+1}$ derived from that estimate are sensitive to our use of equation (23) and data on the current account in our measurement procedure.


% Source PDF 39, text block 3.
We do so for two reasons. First, recall our measurement procedure identifies $r_{t+1}^{*}-\bar{g}_{t+1}$ from the US current account. However, prior papers such as Farhi and Gourio (2018) and Crouzet and Eberly (2023) do not consider the implications of their models for the current account and instead use data on historical growth rates to estimate $\bar{g}_{t+1}$. Second, one might consider an alternative model structure that does not have a representative US household. For example, Greenwald, Lettau, and Ludvigson (2025) consider a model in which there are two types of US households, one that earns labor income and consumes that income every period (living hand to mouth) and another that owns all financial assets. If we were to make a similar assumption, our model’s implication for the current account (equation (23)) would not include the terms involving current labor compensation $W_{t}L_{t}$ and the level of human wealth $H_{t}$ since the households earning labor income would contribute nothing to the aggregate saving rate.


% Source PDF 39, text block 4.
To conduct our sensitivity analysis, we consider four alternative measurement procedures in which we drop equation (23) for the current account and instead introduce an alternative auxiliary assumption about either expected growth $\bar{g}_{t+1}$ or the expected dividend yield $r_{t+1}^{*}-\bar{g}_{t+1}$. Specifically, we calculate the model implied paths for $r_{t+1}^{*},\mu _{t+1}$, and $\alpha _{t+1}$ under the auxiliary assumptions that


% Source PDF 39, text block 5.
(i) $r_{t+1}^{*}-\bar{g}_{t+1}$ is constant at the average of $D_{t+1}/V_{t}$ over the 1990–2023 time period;


% Source PDF 39, text block 6.
(ii) $\bar{g}_{t+1}$ is equal each quarter to the median 10-year GDP growth forecast in the Survey of Professional Forecasters, collected by the Federal Reserve Bank of Philadelphia (2024);\hyperlink{fn:31}{\textsuperscript{31}}


% Source PDF 39, text block 7.
(iii) $\bar{g}_{t+1}$ is set equal to the trend from an HP filtered series of quarterly growth rates of US corporate GVA; and


% Source PDF 39, text block 8.
(iv) the expected dividend yield $r_{t+1}^{*}-\bar{g}_{t+1}$ is equal to the realized dividend yield $D_{t+1}/V_{t}$ each period.


% Source PDF 39, text block 9.
We show the results for $r_{t}^{*},\bar{g}_{t},\mu _{t}$, and $\alpha _{t}$ under these four alternative measurement scenarios in Figure 15. We include in this plot (in green) the values of these parameters that we obtain in our baseline measurement exercise in which our model with a representative US household replicates the path for the current account.


% Source PDF 39, text block 10.
Note, first, that the series for $r_{t}^{*},\bar{g}_{t},\mu _{t}$, and $\alpha _{t}$ that we obtain under the assumption that $r_{t+1}^{*}-\bar{g}_{t+1}$ is constant at the average of $D_{t+1}/V_{t}$ over the 1990–2023 time period (blue) are so close to their counterparts in the baseline calibration (green) that it is


% Source PDF 39, text block 11.
{\footnotesize \hypertarget{fn:31}{}\textbf{Footnote 31.} We linearly interpolate the original data to develop a quarterly series for $\bar{g}_{t+1}$.\par}



% ===== ORIGINAL PDF PAGE 40; JOURNAL PAGE 2190 =====
\sourcepage{40}{2190}

% Figure 15: original PDF page 40, plot crop x=60..444, y=62..367 points.
\begin{center}
\begingroup\renewcommand{\thefigure}{15}\refstepcounter{figure}\label{fig:15}\endgroup
\includegraphics[page=40,trim=60 353 60 62,clip,width=0.95\linewidth]{\sourcepdf}
\par\small Figure 15. Sensitivity to Trend Growth Model
\end{center}


% Source PDF 40, text block 1.
difficult to distinguish the two alternative measurements in the figure. The reason for this similarity is that $r_{t+1}^{*}-\bar{g}_{t+1}$ is close to constant in our baseline calibration. The time paths for $\bar{g}_{t+1}$ under the other specifications vary widely. Nonetheless, these alternative measurement procedures produce time series for other parameter values that are similar to those in our baseline specification, outside of the years around the peak of the dot-com boom in 2000 and the stock market boom in the last few years of our sample.


% Source PDF 40, text block 2.
To understand this, note that when Tobin’s Q is equal to one, the calibration procedure will deliver an output wedge $\mu _{t+1}$ equal to one, independent of the values for $r_{t+1}^{*}$ and $\bar{g}_{t+1}$. In that case, $r_{t+1}^{*}$ and $\alpha _{t+1}$ are identified solely from labor’s share and from the Euler equation for investment (equations (9) and (12)), and because neither of those equations involves future trend growth, the imputed values for $r_{t+1}^{*}$ and $\alpha _{t+1}$ are independent of $\bar{g}_{t+1}$. That is why the paths for $r_{t+1}^{*}$ and $\alpha _{t+1}$ are similar across alternative models for expected growth in periods when Q (and thus $\mu _{t+1}$) is close to one. The further we move away from Q  $=$ 1, the larger is the impact of $\bar{g}_{t+1}$ on the imputed value for $\mu _{t+1}$, which translates into larger impacts on the estimated values for $\alpha _{t+1}$ and $r_{t+1}^{*}$. For example, the SPF growth forecast in 2000 (the red line in the top right panel of Figure 15) was quite strong compared to our baseline parameterization. Thus, that model requires a smaller output wedge $\mu _{t+1}$ to match observed valuations in 2000, which in turn calls for a higher value for $\alpha _{t+1}$ to match labor’s



% ===== ORIGINAL PDF PAGE 41; JOURNAL PAGE 2191 =====
\sourcepage{41}{2191}

% Source PDF 41, text block 1.
share. Because a higher $\alpha _{t+1}$ implies higher rental income from capital, the imputed cost of capital $r_{t+1}^{*}$ is also higher. While different models for trend growth $\bar{g}_{t+1}$ deliver generally similar estimates for other model parameters, they generate wildly different predictions for the dynamics of the US current account and thus the US net foreign asset position. The reason is that the value of human wealth and thus desired consumption are very sensitive to the valuation multiple $(r_{t+1}^{*}-\bar{g}_{t+1})^{-1}$. Thus, in the context of our model, the baseline path for $\bar{g}_{t+1}$ is strongly preferred to the alternatives we have considered.


% Source PDF 41, text block 2.
\section*{VII. Conclusion}
\addcontentsline{toc}{section}{VII. Conclusion}


% Source PDF 41, text block 3.
Cross-border asset holdings have grown very large in recent decades. These large gross positions open the door to new channels for shocks to propagate internationally, especially shocks that affect asset values. Gourinchas and Rey (2014) showed that changes in asset values play a quantitatively important role, in an accounting sense, in driving the dynamics of the NFA position. We show that that finding extends with even more force in recent data: The unprecedented decline in US NFA position since the Great Recession is primarily driven by a boom in the value of corporate equity that foreigners own in the United States.


% Source PDF 41, text block 4.
We have presented a simple international macrofinance model to measure and interpret the factors driving changes in value of the US corporate sector and the US current account and NFA position over the period 1990–2023. Our model extends previous macrofinance models used for similar measurement exercises in integrating the evolution of the current account. In doing so, we reinforce previous findings that increases in the cash flows to firm owners, rather than changes in the valuation multiple of those cash flows, account for much of the increase in the value of the US corporate sector over the past decade. Indeed, the share of US corporate GVA available as free cash flow to owners of corporations has reached levels not previously seen in post-WWII data.


% Source PDF 41, text block 5.
Our model extends previous work in international macroeconomics in developing an accounting of observed changes in the US current account and NFA position that incorporates quantitatively realistic fluctuations in asset values. We find that the direct impact of changes in the valuation of the US corporate sector on the US NFA position through its mechanical impact on the valuation of ROW equity holdings in the US has been large over the past decade, while the indirect impact of these developments on the current account through induced changes in the wealth to income ratio for US households has been quite small.


% Source PDF 41, text block 6.
Through the lens of our model, a rising share of factorless income in the United States is a key driver of rising US asset valuations. This rise would not have mattered much for US households absent foreign ownership of US equity. But given high observed ROW ownership of US firms, the rise in US equity values during the 2010s was associated with a large increase in the portion of free cash from US firms accruing to foreign owners, and a large consumption loss for American households.


% Source PDF 41, text block 7.
While our model provides a transparent framework for interpreting macro and financial data, maintaining tractability required assuming risk-neutral foreign investors and an exogenous world interest rate. Modeling risk and risk aversion explicitly would introduce asset-specific risk premia, generating an additional possible driver



% ===== ORIGINAL PDF PAGE 42; JOURNAL PAGE 2192 =====
\sourcepage{42}{2192}

% Source PDF 42, text block 1.
of valuations. We make some progress in that direction in Atkeson, Heathcote, and Perri (2025). Relatedly, in a fully general equilibrium multicountry setting, shocks in all countries would drive the current account, valuation effects, and NFA dynamics. Constructing and estimating such a model is an important avenue for future research.


% Source PDF 42, text block 2.
\section*{Appendix A. Data Sources}
\addcontentsline{toc}{section}{Appendix A. Data Sources}


% Source PDF 42, text block 3.
Our main sources of US data are quarterly and annual versions of the following tables, all contained in the Financial Accounts of the United States (Board of Governors of the Federal Reserve System 2024).


% Source PDF 42, text block 4.
\textbullet{} Table S.1 Selected Aggregates for Total Economy and Sectors of the Integrated Macroeconomic Accounts \textbullet{} Table  S.5 Non Financial Corporate Business sector of the Integrated Macroeconomic Accounts \textbullet{} Table S.6 Financial Business sector of the Integrated Macroeconomic Accounts \textbullet{} Table S.9 Rest of World sector of the Integrated Macroeconomic Accounts \textbullet{} Table B.1 Derivation of US Net Wealth Financial Accounts of the United States


% Source PDF 42, text block 5.
Our main sources for international data are the OECD Annual Nonfinancial Accounts by institutional sector (Expenditure) and the OECD Annual Financial Balance Sheets (stocks), nonconsolidated. We first describe the series we use to measure the levels of the gross and net foreign asset position for the United States and decomposition of changes in those positions into flows and revaluation effects. We then describe our measures of flows and valuations of the corporate sector. Next, we describe how we measure the extent of foreign ownership of US equities, including the equity for foreign parent firms in their US subsidiaries. Finally, we describe how we construct measures of flows and valuations of the corporate sector for the European Union and the G6. All the tables used to construct the series in the paper and the details of the construction of every series are collected in a single Excel file, which is available on the AEA Data and Code Repository.


% Source PDF 42, text block 6.
\subsection*{A1. Gross and Net Foreign Assets, Flows, and Valuations}
\addcontentsline{toc}{subsection}{A1. Gross and Net Foreign Assets, Flows, and Valuations}


% Source PDF 42, text block 7.
Gross and Net Foreign Assets.—Data on gross and net foreign assets are taken from Table S.9. Table S.9 is presented from the perspective of the rest of the world. We consider flows and net foreign assets from the perspective of the United States. Thus, we typically take the negative of the series noted below. The total market value of financial claims of the United States on the ROW is given in line 136 of Table S.9 (Total liabilities). The total market value of financial claims of the ROW on the United States is given on line 107 (Total financial assets) of Table S.9. These two series constitute the gross foreign asset positions used in our study, with the NFA position of the United States shown in Figure 1 being the difference between the market value of US claims on the ROW and ROW claims on the United States, which corresponds to (the negative of) line 161 (Net worth external account) in Table S.9.



% ===== ORIGINAL PDF PAGE 43; JOURNAL PAGE 2193 =====
\sourcepage{43}{2193}

% Source PDF 43, text block 1.
We take ratios of these and subsequent series relative to nominal GDP (Line 1 from Bureau of Economic Analysis table 1.1.5) and to nominal Gross Value Added of the Corporate Sector, which we construct using quarterly data from Tables S.5 and S.6, as described below.


% Source PDF 43, text block 2.
The Current Account, the Capital Account, and Valuation Changes.—Using data from Table S.9, we decompose nominal changes in the US net foreign asset position according to the following accounting identity:


% Source PDF 43, text block 3.
\[
NFA_t-NFA_{t-1}=\underbrace{CA_t}_{\text{net lending abroad}}+\underbrace{VA_t}_{\text{valuation changes}}+\underbrace{RES_t}_{\text{residual term}}.
\]


% Source PDF 43, text block 16.
The variables $NFA_{t-1}$ and $NFA_{t}$ are the end-of-previous-period and end-of-current-period net foreign asset positions of the United States described above. The change $NFA_{t}-NFA_{t-1}$ is reported (with the opposite sign) on line 106 (Change in Net Worth). The current account $CA_{t}$ measured from the goods and services flow side is the negative of line 14 (Net lending or borrowing). Note that this series is annualized, so we divide the quarterly data by four. “Valuation changes” $VA_{t}$ is the negative of line 105 (Changes in net worth due to nominal holding gains/losses).


% Source PDF 43, text block 17.
The term $RES_{t}$ is reported in line 71 (Total other volume changes) and is equal to line 72 (other volume changes) minus line 73 (statistical discrepancy), which measures the difference between net lending abroad measured from the goods and services flow side and from observed net financial flows, so that


% Source PDF 43, text block 18.
\[
\underbrace{RES_t}_{\text{residual term}}=\underbrace{OV_t}_{\text{other volume changes}}-\underbrace{SD_t}_{\text{statistical discrepancy}}
\]


% Source PDF 43, text block 31.
and
\[
\underbrace{SD_t}_{\text{statistical discrepancy}}=\underbrace{CA_t}_{\text{net lending abroad}}-\underbrace{NFT_t}_{\substack{\text{net lending abroad from}\\\text{financial transactions}}}.
\]


% Source PDF 43, text block 44.
$NFT_{t}$ is reported in line 70 (Net Lending or Borrowing, Financial Account). Note that line 70 is annualized in quarterly data, just like line 14, so we divide it by four. Using the equations above, we can write an alternative decomposition of the changes in the US net foreign asset position as


% Source PDF 43, text block 45.
\[
NFA_t-NFA_{t-1}=\underbrace{NFT_t}_{\text{net financial transactions}}+\underbrace{VA_t}_{\text{valuation changes}}+\underbrace{OV_t}_{\text{other volume changes}}.
\]


% Source PDF 43, text block 58.
As discussed in Bertaut and  Judson (2022), the series for “other volume changes” represents primarily discrepancies arising for separate data sources on gross cross-border asset positions and flows. In Figure A1 we expand the decomposition in Figure 2 in the paper to include the cumulated net financial transactions and the cumulated other volume changes. Note that all these decompositions are invariant to measurement issues within the current account, relating to the measurement of US exports and factor income, as discussed in Guvenen et al. (2022). The figure shows that cumulated other volume changes $(\sum _{j=1}^{t}OV_{j})$ account for very little movement in the net foreign asset position, so that the upward movement in



% ===== ORIGINAL PDF PAGE 44; JOURNAL PAGE 2194 =====
\sourcepage{44}{2194}

% Figure A1: original PDF page 44, plot crop x=60..444, y=62..289 points.
\begin{center}
\begingroup\renewcommand{\thefigure}{A1}\refstepcounter{figure}\label{fig:A1}\endgroup
\includegraphics[page=44,trim=60 431 60 62,clip,width=0.95\linewidth]{\sourcepdf}
\par\small Figure A1. Decomposition of Changes in US Net Foreign Assets over US Corporate Value Added
\end{center}


% Source PDF 44, text block 1.
$\sum _{j=1}^{t}RES_{j}$ term is mostly explained by the statistical discrepancy between $NFT_{t}$ and $CA_{t}$. Specifically, the cumulated net lending computed using financial transactions $(\sum _{j=1}^{t}NFT_{j})$ declines less than the cumulated net lending using trade $(\sum _{j=1}^{t}CA_{j})$, suggesting that using this alternative measure of the net lending increases the importance of negative valuations in accounting the decline in the US net foreign asset position over our sample.


% Source PDF 44, text block 2.
The Equity Component of Gross and Net Foreign Assets.—We measure the equity component of gross and net foreign assets of the United States using the sum of portfolio investments in equity and the equity component of foreign direct investment. The market value of US equity investment in the ROW is given on line 153 of Table S.9, “Equity and Investment Fund Shares.” The market value of ROW equity investment in the United States is reported in line 125 of Table S.9, “Equity and Investment Fund Shares.” We compute market values of non-equity gross and net foreign assets and liabilities as the difference between the measures of the total positions and the equity component of those positions as described above. The corresponding valuation changes of the market valuations of the equity assets and liabilities are reported in the “Revaluation account” section of Table S.9. Specifically, the revaluation of US equity investment in the ROW is given on line 100 of Table S.9, “Equity and investment fund shares.” The revaluation of ROW equity investment in the United States is reported in line 84, “Equity and investment fund shares.”


% Source PDF 44, text block 3.
\subsection*{A2. Measurement of the US Corporate Sector}
\addcontentsline{toc}{subsection}{A2. Measurement of the US Corporate Sector}


% Source PDF 44, text block 4.
We now detail exactly which series we use for each entry.


% Source PDF 44, text block 5.
Gross Value Added.—Gross value added for the nonfinancial corporate business sector is given in line 1 of Table S.5, and that for the financial business sector on



% ===== ORIGINAL PDF PAGE 45; JOURNAL PAGE 2195 =====
\sourcepage{45}{2195}

% Figure A2: original PDF page 45, plot crop x=60..444, y=62..293 points.
\begin{center}
\begingroup\renewcommand{\thefigure}{A2}\refstepcounter{figure}\label{fig:A2}\endgroup
\includegraphics[page=45,trim=60 427 60 62,clip,width=0.95\linewidth]{\sourcepdf}
\par\small Figure A2. US Corporate Sector Gross Value Added Share of GDP
\end{center}


% Source PDF 45, text block 1.
line 1 of Table S.6. We compute the fraction of gross value added in the corporate sector as the sum of that in the nonfinancial corporate business sector and in the financial business sector, all divided by GDP (US Bureau of Economic Analysis (2024b), table 1.1.5). In Figure A2, we show the share of economy-wide gross value added that is produced in the US corporate sector.


% Source PDF 45, text block 2.
The use of the residence principle has a substantial impact on the measurement of economic activity in the corporate sector, relative to what one would get if one were to instead associate the economic activity of affiliates of multinational enterprises with the country in which the multinational is headquartered. For example, the Bureau of Economic Analysis reports that in 2018, majority-owned US affiliates of foreign multinational enterprises contributed \$1.1 trillion, or 7.1 percent of US business sector value added, and accounted for 6.0 percent of total private industry employment in the United States. Likewise, in 2018, US multinationals produced \$5.7 trillion of value added, \$4.2 trillion of which was produced by US resident operations with 28.6 million employees, and \$1.5 trillion of which was produced by majority-owned affiliates abroad with 14.4 million employees. Using the residence principle, the Integrated Macroeconomic Accounts include the \$1.1 trillion of value added by US affiliates of foreign multinationals as a flow attributed to the US corporate sector and do not include the \$1.5 trillion produced by foreign affiliates of US multinational enterprises in this category.


% Source PDF 45, text block 3.
Dividends.—The variable in the model is $D_{t}$, which is a comprehensive measure of payouts to investors in the corporate sector from operations. For the nonfinancial corporate business sector, we measure payouts using the following lines from Table S.5. We take operating surplus, net in line 8 less current taxes on income, wealth, line 13 less net capital formation in line 20. For the financial business sector,



% ===== ORIGINAL PDF PAGE 46; JOURNAL PAGE 2196 =====
\sourcepage{46}{2196}

% Source PDF 46, text block 1.
we measure payouts from the corresponding lines in Table S.6. We take operating surplus, net in line 8 less current taxes on income, wealth, line 11 less capital formation, net in line 18.


% Source PDF 46, text block 2.
Earnings.—The variable $E_{t}$ in the model is a comprehensive measure of the operating earnings of the US corporate sector. In the model $E_{t}=D_{t}+I_{t}-\delta K_{t}$. We construct this measure using our constructed measure of dividends above, adjusted using the following series from Tables S.5 and S.6. For both the financial and the nonfinancial corporate business sector, we add net capital formation to our measure of dividends $D_{t}$.


% Source PDF 46, text block 3.
Replacement Value of Nonfinancial Assets.—The variable $Q_{t}K_{t+1}$ in the model is the replacement value of nonfinancial assets at the end of period $t$. This is the sum of such values across the nonfinancial business sector and the financial business sector. We construct this measure as the sum of line 102 (Non Financial assets) in Table S.5 and line 97 (Non Financial Assets) in Table S.6. These include structures, equipment, intellectual property products, and inventories.


% Source PDF 46, text block 4.
Market or Enterprise Value of Corporate Nonfinancial Assets.—The variable $V_{t}$ in the model is the market or enterprise value of nonfinancial assets at the end of period $t$. This is the sum of such values across the nonfinancial corporate sector and the financial business sector. These measures are available in Table B.1, as market value of the nonfinancial assets. Specifically, for the nonfinancial corporate sector, we use line 14 (Nonfinancial corporate business; market value estimate of nonfinancial assets) and for the financial corporate sector, we use line 15 (Domestic financial sectors; market value estimate of nonfinancial assets). As reported in the table, both values are constructed in similar fashion as the market value of corporate equity, plus foreign direct investment: equity, plus miscellaneous other equity (excluding proprietors’ equity), plus total liabilities, less total financial assets.



% ===== ORIGINAL PDF PAGE 47; JOURNAL PAGE 2197 =====
\sourcepage{47}{2197}

% Source PDF 47, text block 1.
\subsection*{A3. Additional Measures of Firm Value and Cash Flow}
\addcontentsline{toc}{subsection}{A3. Additional Measures of Firm Value and Cash Flow}


% Figure A3: original PDF page 47, plot crop x=60..444, y=94..324 points.
\begin{center}
\begingroup\renewcommand{\thefigure}{A3}\refstepcounter{figure}\label{fig:A3}\endgroup
\includegraphics[page=47,trim=60 396 60 94,clip,width=0.95\linewidth]{\sourcepdf}
\par\small Figure A3. Enterprise Value and Replacement Value\par Notes: Red: Total Corporate Sector relative to Total Corporate Gross Value Added. Blue: Nonfinancial Corporate Sector relative to Nonfinancial Corporate Sector Gross Value Added.
\end{center}


% Figure A4: original PDF page 47, plot crop x=60..444, y=390..620 points.
\begin{center}
\begingroup\renewcommand{\thefigure}{A4}\refstepcounter{figure}\label{fig:A4}\endgroup
\includegraphics[page=47,trim=60 100 60 390,clip,width=0.95\linewidth]{\sourcepdf}
\par\small Figure A4. Free Cash Flow: Total Corporate Sector Relative to Total Corporate Gross Value Added and Nonfinancial Corporate Sector Relative to Nonfinancial Corporate Sector Gross Value Added
\end{center}



% ===== ORIGINAL PDF PAGE 48; JOURNAL PAGE 2198 =====
\sourcepage{48}{2198}

% Figure A5: original PDF page 48, plot crop x=60..444, y=62..291 points.
\begin{center}
\begingroup\renewcommand{\thefigure}{A5}\refstepcounter{figure}\label{fig:A5}\endgroup
\includegraphics[page=48,trim=60 429 60 62,clip,width=0.95\linewidth]{\sourcepdf}
\par\small Figure A5. Nonfinancial Corporate Enterprise Value, Nonfinancial Corporate Market Value of Equity, and Nonfinancial Corporate Market Value of Publicly Traded Equity, All Relative to Nonfinancial Corporate Gross Value Added
\end{center}


% Source PDF 48, text block 1.
\subsection*{A4. Foreign Ownership of US Equity and US Ownership of Foreign Equity}
\addcontentsline{toc}{subsection}{A4. Foreign Ownership of US Equity and US Ownership of Foreign Equity}


% Source PDF 48, text block 2.
Our measure of the share of US equity owned by the ROW at the end of period $t,(1-\lambda _{t},)$ is the ratio between the market value of ROW equity assets in the United States, as defined above, and the enterprise value of corporate nonfinancial assets, as defined above. Figure A6 plots the evolution of this ratio.


% Source PDF 48, text block 3.
We do not construct a direct measure of the enterprise value of ROW corporations $V_{t}^{*}$. Instead, we measure the value of US ownership of foreign equity assets at the end of period $t,(\lambda _{t}^{*}V_{t}^{*})$ directly as the market value US equity assets in ROW, as defined above. We measure the dividends that US residents receive in period $t$ on their ownership of equity in the ROW (denoted by $\lambda _{t-1}^{*}D_{t}^{*}$ in the model) as monetary dividends paid, as reported in line 11 from table 4.1 in Bureau of Economic Analysis (2024b). Note, as we discuss below, this series includes only the monetary dividends actually paid on US FDI in the ROW as opposed to the full accounting income reported as part of the current account.


% Source PDF 48, text block 4.
We note that, in contrast to the well-known “income puzzle” on US FDI in the ROW (discussed below), these dividends on US equity in the ROW actually paid are not high relative to the market valuation of US residents holdings of equity in the ROW. We compute that dividend yield as follows. To compute $\lambda _{t-1}^{*}V_{t}^{*}$, we take the market value of US equity in the ROW at the end of period $t-1(\lambda _{t-1}^{*}V_{t-1}^{*})$ as measured above and add to it the revaluation of US equity in the ROW in period $t(\lambda _{t-1}^{*}(V_{t}^{*}-V_{t-1}^{*}))$. The revaluation of US equity investment in the ROW is reported on line 100 of Table S.9.


% Source PDF 48, text block 5.
With these data, we compute the current dividend yield on US equity in ROW as the ratio of $\lambda _{t-1}^{*}D_{t}^{*}$ to $\lambda _{t-1}^{*}V_{t}^{*}$ constructed as above giving $D_{t}^{*}/V_{t}^{*}$, which is plotted



% ===== ORIGINAL PDF PAGE 49; JOURNAL PAGE 2199 =====
\sourcepage{49}{2199}

% Figure A6: original PDF page 49, plot crop x=60..444, y=62..289 points.
\begin{center}
\begingroup\renewcommand{\thefigure}{A6}\refstepcounter{figure}\label{fig:A6}\endgroup
\includegraphics[page=49,trim=60 431 60 62,clip,width=0.95\linewidth]{\sourcepdf}
\par\small Figure A6. ROW Ownership Share of US Corporations, $(1-\lambda _{t})$
\end{center}


% Source PDF 49, text block 1.
in Figure A7. While it is the case that this ratio shows big spikes around changes in US tax law that encourage the repatriation of earnings on US FDI, the long-term average for this series is in line with the current dividend yield on US corporations reported in the right panel of Figure 8.


% Source PDF 49, text block 2.
\subsection*{A5. Measurement of the Foreign Corporate Sector}
\addcontentsline{toc}{subsection}{A5. Measurement of the Foreign Corporate Sector}


% Source PDF 49, text block 3.
The series in Figure 9 are computed as follows. The series for enterprise value, cash flows, and gross value added for the corporate sector for the Unites States are computed as described above, except for the fact that the series in these figures are annual instead of quarterly. For the European Union and each single country in the G6 (Canada, France, Germany, Italy, Japan, and the United Kingdom), we compute the enterprise value of the corporate sector as the sum of the “net financial worth” of nonfinancial and of financial corporations, all from the OECD “Annual Financial Balance Sheets (stocks), nonconsolidated.” The gross value added of the corporate sector is from the OECD “Annual non-financial accounts by institutional sector (Expenditure),” while the cash flows are computed as “Operating surplus and mixed income, gross” minus “Gross capital formation” and minus “Current taxes on income, wealth, etc.” from the same OECD dataset. Note that for Japan and Canada, the OECD does not report gross value added for the corporate sector, so for those two countries, we estimate a share of the corporate sector by taking the ratio of “Operating surplus and mixed income” of the corporate sector to the same variable for the whole economy and then use the share multiplied by gross value added of the whole economy to obtain a series for gross value added of the corporate sector. We also check that when we construct figures for gross value added, enterprise value, and cash flows for the United States using OECD datasets, the series are comparable to those that we construct using US data sources. The only discrepancy is that the



% ===== ORIGINAL PDF PAGE 50; JOURNAL PAGE 2200 =====
\sourcepage{50}{2200}

% Figure A7: original PDF page 50, plot crop x=60..444, y=62..288 points.
\begin{center}
\begingroup\renewcommand{\thefigure}{A7}\refstepcounter{figure}\label{fig:A7}\endgroup
\includegraphics[page=50,trim=60 432 60 62,clip,width=0.95\linewidth]{\sourcepdf}
\par\small Figure A7. Current Dividend Yield on US Equity in ROW, $D_{t}^{*}/V_{t}^{*}$
\end{center}


% Source PDF 50, text block 1.
net worth of nonfinancial corporations for the United States reported by the OECD include the net worth of nonfinancial noncorporate businesses. Series for G6 countries are computed aggregating individual countries figures using (period average) market exchange rates. We have also experimented using PPP exchange rates, and results are quantitatively very similar. Exchange rate figures are from the OECD dataset “Annual Purchasing Power Parities and exchange rates.” All OECD datasets are available on OECD (2024).


% Source PDF 50, text block 2.
\section*{Appendix B. Ex Ante Risk Sharing}
\addcontentsline{toc}{section}{Appendix B. Ex Ante Risk Sharing}


% Source PDF 50, text block 3.
In this paper, we have focused on the ex post welfare implications of the recent boom in the value of US Corporations and how those ex post welfare implications are impacted by the extent of gross cross-border equity positions. There is a large literature in international macroeconomics that considers ex ante optimal cross-border equity holdings and how the extent of ex ante optimal international diversification depends on the nature of the shocks hitting the economy and the structure of international trade in goods. See, for example, Baxter and Jermann (1997) and Heathcote and Perri (2013). Coeurdacier, Kollmann, and Martin (2007) is perhaps most relevant for our study in that they consider shocks that reallocate the distribution of income between workers and owners of firms. One main result in their study is that, in the face of such shocks, rather limited cross-border equity positions are ex ante optimal. We next derive that result in our model.


% Source PDF 50, text block 4.
\subsection*{B1. Proofs of Propositions 1 and 2}
\addcontentsline{toc}{subsection}{B1. Proofs of Propositions 1 and 2}


% Source PDF 50, text block 5.
We start with Proposition 1, the scenario in which there are news shocks about the future profit share.



% ===== ORIGINAL PDF PAGE 51; JOURNAL PAGE 2201 =====
\sourcepage{51}{2201}

% Source PDF 51, text block 1.
Given logarithmic utility, the domestic household’s first-order condition for risk-free saving in the globally traded bond is


% Source PDF 51, text block 2.
\[
\frac{1}{C_t}=\frac{1}{(1+\rho)}(1+r^*)\mathbb{E}_t\left[\frac{1}{C_{t+1}}\right].
\]


% Source PDF 51, text block 6.
If $C_{t+1}=(1+\bar g)C_t$ and $\frac{1}{1+\rho}\frac{1}{1+\bar g}=\frac{1}{1+r^*}$, this condition is satisfied at each date $t$.


% Source PDF 51, text block 9.
Imposing the conjectured consumption rule, the domestic household’s FOC for risky domestic equity purchases is


% Source PDF 51, text block 10.
\[
\frac{V_t}{C_t}=\frac{1}{(1+\rho)}\frac{1}{(1+\bar g)C_t}\mathbb{E}_t[V_{t+1}+D_{t+1}].
\]


% Source PDF 51, text block 14.
Recall that the foreign representative household is risk neutral. Thus, all financial assets in the economy, including domestic equity, must offer the same expected return $r^{*}$, implying


% Source PDF 51, text block 15.
\[
\frac{\mathbb{E}_t[V_{t+1}+D_{t+1}]}{V_t}=1+r^*.
\]


% Source PDF 51, text block 16.
It follows that the domestic household’s FOC for domestic equity purchases is satisfied.


% Source PDF 51, text block 17.
The only condition that remains to be checked is that the conjectured allocation is budget feasible for domestic households.


% Source PDF 51, text block 18.
To verify this, we will solve for equilibrium date $0$ consumption, $C_{0}$, and show that given $C_{0}$, the present value of lifetime wealth will grow at the constant rate $\bar{g}$ when consumption grows at rate $\bar{g}$, independent of the sequence of realizations for the profit-share news shock.


% Source PDF 51, text block 19.
Let $\mathit{Wealth}_{t}$ denote the present value of the representative domestic household’s income plus end-of-period wealth at date $t$:


% Source PDF 51, text block 20.
\[
\begin{aligned}
\mathit{Wealth}_t={}&W_tL_t+\lambda_{t-1}(D_t+V_t)+\lambda_{t-1}^*(D_t^*+V_t^*)\\
&+(1+r^*)B_t+H_t,
\end{aligned}
\]


% Source PDF 51, text block 21.
where $H_{t}$ is the present value of human wealth.


% Source PDF 51, text block 22.
The value of human wealth is the expected present value of future labor income. At any date $t$, expected labor income at $t+1$ is $(1-\kappa_{t+1})(1-\alpha)(1-\tau)(1+\bar g)Y_t$. From $t+1$ onward, $\kappa_t$ follows a unit root process such that $\mathbb{E}_t[\kappa_{t+j}]=\kappa_{t+1}$, while income $Y_{t+j}$ will grow at constant rate $\bar g$. Given the conjectured equilibrium process for consumption, according to which consumption growth is independent of shocks to the share of factorless income, labor income in all future states is discounted at rate $\frac{1}{1+\rho}\frac{1}{1+\bar g}=\frac{1}{1+r^*}$. Thus,


% Source PDF 51, text block 25.
\[
\begin{aligned}
H_t&=\mathbb{E}_t\sum_{j=1}^{\infty}\frac{(1-\kappa_{t+j})(1-\alpha)(1-\tau)Y_{t+j}}{(1+r^*)^j}\\
&=\frac{(1-\kappa_{t+1})(1-\alpha)(1-\tau)(1+\bar g)Y_t}{r^*-\bar g}.
\end{aligned}
\]


% Source PDF 51, text block 30.
It is useful to rewrite the expression for lifetime wealth in a slightly different way, conceptualizing domestic firms as comprising two pieces: (i) ownership of capital,



% ===== ORIGINAL PDF PAGE 52; JOURNAL PAGE 2202 =====
\sourcepage{52}{2202}

% Source PDF 52, text block 1.
with value $V_{t}^{K}=QK_{t+1}$, and (ii) ownership of future claims to monopoly profits, with value


% Source PDF 52, text block 2.
\[
V_t^{\Pi}=\mathbb{E}_t\sum_{j=1}^{\infty}\frac{\kappa_{t+j}(1-\tau)Y_{t+j}}{(1+r^*)^j}=\frac{\kappa_{t+1}(1-\tau)(1+\bar g)Y_t}{r^*-\bar g}.
\]


% Source PDF 52, text block 6.
News shocks have no impact on the return to the physical capital component of firm value because at each date $t$ when news about $\kappa _{t+1}$ is revealed, investment can be adjusted so that at $t+1$, the rental rate of capital net of depreciation, $R_{t+1}-\delta $, will be equal to $r^{*}$. Thus, the only component of financial wealth whose return is risky is the claim-to-profits portion of domestic equity. We can therefore express lifetime wealth at $t$ as


% Source PDF 52, text block 7.
\[
\mathit{Wealth}_t=W_tL_t+\lambda_{t-1}(\Pi_t+V_t^{\Pi})+(1+r^*)\widetilde B_t+H_t,
\]


% Source PDF 52, text block 8.
where $\tilde{B}_{t}$ lumps together all the safe components of wealth:


% Source PDF 52, text block 9.
\[
\begin{aligned}
(1+r^*)\widetilde B_t={}&\lambda_{t-1}[R_tK_t+(1-\delta)QK_t-QK_{t+1}]+\lambda_{t-1}QK_{t+1}\\
&+\lambda_{t-1}^*(D_t^*+V_t^*)+(1+r^*)B_t\\
={}&(1+r^*)\left(B_t+\lambda_{t-1}QK_{t-1}+\lambda_{t-1}^*V_{t-1}^*\right).
\end{aligned}
\]% Conversion note: the last line prints QK_{t-1}; preserved as printed.


% Source PDF 52, text block 12.
At date $0,\tilde{B}_{0}$ is given by the expression above given the initial state variables $\lambda _{-1}QK_{t-1},\lambda _{-1}^{*}V_{t-1}^{*}$ and $B_{0}$.

Given log utility, optimal consumption at each date $t$ is


% Source PDF 52, text block 13.
\[
C_t=\frac{\rho}{1+\rho}\mathit{Wealth}_t.
\]


% Source PDF 52, text block 14.
Given $C_0=\frac{\rho}{1+\rho}\mathit{Wealth}_0$, $\widetilde B_1$ is given by the sequential budget constraint,


% Source PDF 52, text block 15.
\[
C_0+(\lambda_0-\lambda_{-1})V_0^{\Pi}+\widetilde B_1=W_0L_0+\lambda_{-1}\Pi_0+(1+r^*)\widetilde B_0,
\]


% Source PDF 52, text block 16.
with $\lambda _{0}=1-\alpha $. Thereafter, there is no further trade in domestic equity.


% Source PDF 52, text block 17.
It remains to check that, given the rule for consumption $C_t=\frac{\rho}{1+\rho}\mathit{Wealth}_t$ and the equity rule $\lambda_t=(1-\alpha)$, total wealth and consumption do in fact both grow at rate $1+\bar g$ at every date $t$. Consumption at any date $t$ is


% Source PDF 52, text block 18.
\[
\begin{aligned}
C_t&=\frac{\rho}{1+\rho}\mathit{Wealth}_t\\
&=\frac{\rho}{1+\rho}\left(W_tL_t+\lambda_{t-1}\Pi_t+\lambda_{t-1}V_t^{\Pi}+H_0+\lambda_{t-1}^*(1+r^*)\widetilde B_t\right)\\
&=\frac{\rho}{1+\rho}\Biggl[(1+r^*)\widetilde B_t+\left[(1-\kappa_t)(1-\alpha)+\lambda_{t-1}\kappa_t\right](1-\tau)Y_t\\
&\qquad+\frac{\left[(1-\kappa_{t+1})(1-\alpha)+\lambda_{t-1}\kappa_{t+1}\right](1-\tau)(1+\bar g)Y_t}{r^*-\bar g}
+(1+r^*)\widetilde B_t\Biggr].
\end{aligned}
\]
% Conversion note: this published display contains H_0, lambda^*_{t-1} multiplying safe wealth, and two occurrences of (1+r^*) B-tilde_t in the final equality. These apparent inconsistencies are preserved, not silently repaired.




% ===== ORIGINAL PDF PAGE 53; JOURNAL PAGE 2203 =====
\sourcepage{53}{2203}

% Source PDF 53, text block 1.
Imposing $\lambda _{t-1}=(1-\alpha )$, this simplifies to


% Source PDF 53, text block 2.
\[
C_t=\frac{\rho}{1+\rho}\left[\frac{(1-\alpha)(1-\tau)(1+r^*)Y_t}{r^*-\bar g}+(1+r^*)\widetilde B_t\right].
\]


% Source PDF 53, text block 3.
Now, $1=\frac{1}{1+\rho}\frac{1+r^*}{1+\bar g}$ implies $\frac{\rho}{1+\rho}=\frac{r^*-\bar g}{1+r^*}$, and thus,


% Source PDF 53, text block 5.
\[
C_t=(1-\alpha)(1-\tau)Y_t+(r^*-\bar g)\widetilde B_t.
\]


% Source PDF 53, text block 6.
Returning to the sequential budget constraint and substituting in this expression for $C_{t}$ gives


% Source PDF 53, text block 7.
\[
\begin{gathered}
C_t+\widetilde B_{t+1}=W_tL_t+\lambda_{t-1}\Pi_t+(1+r^*)\widetilde B_t\\
(1-\alpha)(1-\tau)Y_t+(r^*-\bar g)\widetilde B_t+\widetilde B_{t+1}=(1-\alpha)(1-\tau)Y_t+(1+r^*)\widetilde B_t\\
\widetilde B_{t+1}=(1+\bar g)\widetilde B_t.
\end{gathered}
\]


% Source PDF 53, text block 10.
Thus, this consumption rule and constant equity portfolio allows the domestic agent to passively accumulate safe financial wealth at constant rate $1+\bar{g}$. This proves that the consumption rule is budget feasible.


% Source PDF 53, text block 11.
The proof for the case in which there are shocks in income growth is identical in form. In that case,


% Source PDF 53, text block 12.
\[
\begin{aligned}
H_t&=\mathbb{E}_t\sum_{j=1}^{\infty}\frac{(1-\kappa)(1-\alpha)(1-\tau)Y_{t+j}}{(1+r^*)^j}\\
&=\frac{(1-\kappa)(1-\alpha)(1-\tau)(1+\bar g)\varepsilon_{Zt}Y_t}{r^*-\bar g},
\end{aligned}
\]


% Source PDF 53, text block 17.
\[
V_t^{\Pi}=\mathbb{E}_t\sum_{j=1}^{\infty}\frac{\kappa(1-\tau)Y_{t+j}}{(1+r^*)^j}=\frac{\kappa(1-\tau)(1+\bar g)\varepsilon_{Zt}Y_t}{r^*-\bar g}.
\]


% Source PDF 53, text block 21.
It remains to check that, given the rule for consumption $C_t=\frac{\rho}{1+\rho}\mathit{Wealth}_t$ and the equity rule $\lambda_t=-(1-\alpha)(1-\kappa)/\kappa$, total wealth and consumption do in fact both grow at rate $1+\bar g$ at every date $t$. Consumption is


% Source PDF 53, text block 22.
\[
\begin{aligned}
C_t&=\frac{\rho}{1+\rho}\mathit{Wealth}_t\\
&=\frac{\rho}{1+\rho}\left[W_tL_t+\lambda_{t-1}\Pi_t+\lambda_{t-1}V_t^{\Pi}+H_0+\lambda_{t-1}^*(1+r^*)\widetilde B_t\right]\\
&=\frac{\rho}{1+\rho}\Biggl(\left[(1-\kappa)(1-\alpha)+\lambda_{t-1}\kappa\right](1-\tau)Y_t\\
&\qquad+\frac{\left[(1-\kappa)(1-\alpha)+\lambda_{t-1}\kappa\right](1-\tau)(1+\bar g)\varepsilon_{Zt}Y_t}{r^*-\bar g}
+(1+r^*)\widetilde B_t\Biggr).
\end{aligned}
\]% Conversion note: H_0 and lambda^*_{t-1} in the second line are preserved as printed.



% ===== ORIGINAL PDF PAGE 54; JOURNAL PAGE 2204 =====
\sourcepage{54}{2204}

% Source PDF 54, text block 1.
Imposing $\lambda _{t-1}=-(1-\alpha )(1-\kappa )/\kappa $, this simplifies to


% Source PDF 54, text block 2.
\[
C_t=\frac{\rho}{1+\rho}(1+r^*)\widetilde B_t.
\]


% Source PDF 54, text block 3.
Now, $1=\frac{1}{1+\rho}\frac{1+r^*}{1+\bar g}$ implies $\frac{\rho}{1+\rho}=\frac{r^*-\bar g}{1+r^*}$, and thus,


% Source PDF 54, text block 5.
\[
C_t=(r^*-\bar g)\widetilde B_t.
\]


% Source PDF 54, text block 6.
Returning to the sequential budget constraint and substituting in this expression for $C_{t}$ gives


% Source PDF 54, text block 7.
\[
\begin{gathered}
C_t+\widetilde B_{t+1}=W_tL_t+\lambda_{t-1}\Pi_t+(1+r^*)\widetilde B_t,\\
(r^*-\bar g)\widetilde B_t+\widetilde B_{t+1}=(1+r^*)\widetilde B_t,\\
\widetilde B_{t+1}=(1+\bar g)\widetilde B_t.
\end{gathered}
\]


% Source PDF 54, text block 10.
Thus, this consumption rule and constant equity portfolio again allows the domestic agent to passively accumulate safe financial wealth at constant rate $1+\bar{g}$, proving that the consumption rule is budget feasible.


% Source PDF 54, text block 11.
\section*{References}
\addcontentsline{toc}{section}{References}


% Source PDF 54, text block 12.
% Reference 1; original PDF 54.
\noindent\hangindent=1.5em\hangafter=1 Aguiar, Mark, and Gita Gopinath. 2007. “Emerging Market Business Cycles: The Cycle Is the Trend.” Journal of Political Economy 115 (1): 69–102.\par

% Reference 2; original PDF 54.
\noindent\hangindent=1.5em\hangafter=1 Atkeson, Andrew, and Ariel Burstein. 2007. “Pricing-to-Market in a Ricardian Model of International Trade.” American Economic Review 97 (2): 362–67.\par

% Reference 3; original PDF 54.
\noindent\hangindent=1.5em\hangafter=1 Atkeson, Andrew, Jonathan Heathcote, and Fabrizio Perri. 2025. “Reconciling Macroeconomics and Finance for the US Corporate Sector: 1929 to Present.” NBER Working Paper 33459.\par

% Reference 4; original PDF 54.
\noindent\hangindent=1.5em\hangafter=1 Atkeson, Andrew, Jonathan Heathcote, and Fabrizio Perri. 2025. Data and Code for: “The End of Privilege: A Reexamination of the Net Foreign Asset Position of the United States.” Nashville, TN: American Economic Association; distributed by Inter-university Consortium for Political and Social Research, Ann Arbor, MI. \url{https://doi.org/10.3886/E209797V1}.\par

% Reference 5; original PDF 54.
\noindent\hangindent=1.5em\hangafter=1 Auclert, Adrien, Hannes Malmberg, Frederic Martenet, and Matthew Rognlie. 2021. “Demographics, Wealth, and Global Imbalances in the Twenty-First Century.” NBER Working Paper 29161.\par

% Reference 6; original PDF 54.
\noindent\hangindent=1.5em\hangafter=1 Avdjiev, Stefan, Mary Everett, Phillip R. Lane, and Hyun Song Shin. 2018. “Tracking the International Footprints of Global Firms.” BIS Quarterly Review 19 (1): 47–66.\par

% Reference 7; original PDF 54.
\noindent\hangindent=1.5em\hangafter=1 Baqaee, David Rezza, and Emmanuel Farhi. 2020. “Productivity and Misallocation in General Equilibrium.” Quarterly Journal of Economics 135 (1): 105–63.\par

% Reference 8; original PDF 54.
\noindent\hangindent=1.5em\hangafter=1 Barkai, Simcha. 2020. “Declining Capital and Labor Shares.” Journal of Finance 75 (5): 2421– 63.\par

% Reference 9; original PDF 54.
\noindent\hangindent=1.5em\hangafter=1 Baxter, Marianne, and Urban J. Jermann. 1997. “The International Diversification Puzzle Is Worse Than You Think.” American Economic Review 87 (1): 170–80.\par

% Reference 10; original PDF 54.
\noindent\hangindent=1.5em\hangafter=1 Bernard, Andrew B., Jonthan Eaton, J. Bradford Jensen, and Samuel Kortum. 2003. “Plants and Productivity in International Trade.” American Economic Review 93 (4): 1268–90.\par

% Reference 11; original PDF 54.
\noindent\hangindent=1.5em\hangafter=1 Bertaut, Carol C., Beau Bressler, and Stephanie Curcuru. 2019. “Globalization and the Geography of Capital Flows.” FEDS Notes, September 6. \url{https://doi.org/10.17016/238--7172.2446}.\par

% Reference 12; original PDF 54.
\noindent\hangindent=1.5em\hangafter=1 Bertaut, Carol, Stephanie E. Curcuru, Ester Faia, and Pierre-Olivier Gourinchas. 2024. “The Global (Mis)Allocation of Capital.” International Finance Discussion Paper 1399.\par

% Reference 13; original PDF 54.
\noindent\hangindent=1.5em\hangafter=1 Bertaut, Carol, and Ruth Judson. 2022. “Estimating US Cross-Border Securities Flows: Ten Years of the TIC SLT.” FEDS Notes, February 18. \url{https://doi.org/10.17016/238-07172.3068}.\par

% Reference 14; original PDF 54.
\noindent\hangindent=1.5em\hangafter=1 Board of Governors of the Federal Reserve System. 2024. Financial Accounts of the United States - Z.1. \url{https://www.federalreserve.gov/releases/z1/20240307/html/default.htm}, March 2024 Release.\par

% Reference 15; original PDF 54.
\noindent\hangindent=1.5em\hangafter=1 Bureau of Economic Analysis. 2014. US International Economic Accounts: Concepts and Methods. United States Department of Commerce.\par



% ===== ORIGINAL PDF PAGE 55; JOURNAL PAGE 2205 =====
\sourcepage{55}{2205}

% Source PDF 55, text block 1.
% Reference 16; original PDF 55.
\noindent\hangindent=1.5em\hangafter=1 Bureau of Economic Analysis. 2024a. International Data, International Transactions, International Services, and International Investment Position, Table 1.3. March 2024 Release.\par

% Reference 17; original PDF 55.
\noindent\hangindent=1.5em\hangafter=1 Bureau of Economic Analysis. 2024b. National Data, National Income and Product Accounts, Tables 1.1.4, 1.1.5, 4.1. March 2024 Release.\par

% Reference 18; original PDF 55.
\noindent\hangindent=1.5em\hangafter=1 Caballero, Ricardo J., Emmanuel Farhi, and Pierre-Olivier Gourinchas. 2017. “Rents, Technical Change, and Risk Premia: Accounting for Secular Trends in Interest Rates, Returns on Capital, Earning Yields, and Factor Shares.” American Economic Review 107 (5): 614–20.\par

% Reference 19; original PDF 55.
\noindent\hangindent=1.5em\hangafter=1 Chen, Zefeng, Zhengyang Jiang, Hanno N. Lustig, Stijn Van Nieuwerburgh, and Mindy Z. Xiaolan. 2022. “Exorbitant Privilege Gained and Lost: Fiscal Implications.” \url{http://dx.doi.org/10.2139/ssrn.4096091}.\par

% Reference 20; original PDF 55.
\noindent\hangindent=1.5em\hangafter=1 Choi, Jason, Rishabh Kirpalani, and Diego Perez. Forthcoming. “US Public Debt and Safe Asset Market Power.” Journal of Political Economy.\par

% Reference 21; original PDF 55.
\noindent\hangindent=1.5em\hangafter=1 Coeurdacier, Nicolas, Robert Kollmann, and Philippe Martin. 2007. “International Portfolios with Supply, Demand, and Redistributive Shocks.” NBER International Seminar on Macroeconomics 4: 231–81.\par

% Reference 22; original PDF 55.
\noindent\hangindent=1.5em\hangafter=1 Coeurdacier, Nicolas, and Hélène Rey. 2013. “Home Bias in Open Economy Financial Macroeconomics.” Journal of Economic Literature 51 (1): 63–115.\par

% Reference 23; original PDF 55.
\noindent\hangindent=1.5em\hangafter=1 Coppola, Antonio, Matteo Maggiori, Brent Neiman, and Jesse Schreger. 2021. “Redrawing the Map of Global Capital Flows: The Role of Cross-Border Financing and Tax Havens.” Quarterly Journal of Economics 136 (3): 1499–1556.\par

% Reference 24; original PDF 55.
\noindent\hangindent=1.5em\hangafter=1 Corhay, Alexandre, Howard Kung, and Lukas Schmid. 2025. “Q: Risks, Rents, or Growth?” Journal of Financial Economics 165: 103990.\par

% Reference 25; original PDF 55.
\noindent\hangindent=1.5em\hangafter=1 Corsetti, Giancarlo, and Panagiotis T. Konstantinou. 2012. “What Drives US Foreign Borrowing? Evidence on the External Adjustment to Transitory and Permanent Shocks.” American Economic Review 102 (2): 1062–92.\par

% Reference 26; original PDF 55.
\noindent\hangindent=1.5em\hangafter=1 Crouzet, Nicolas, and Janice Eberly. 2018. “Understanding Weak Capital Investment: The Role of Market Concentration and Intagibles.” In Changing Market Structures and Implications for Monetary Policy: A Symposium Sponsored by the Federal Reserve Bank of Kansas City, Jackson Hole, WY, August 23–25, 2018, 87–148. Federal Reserve Bank of Kansas City.\par

% Reference 27; original PDF 55.
\noindent\hangindent=1.5em\hangafter=1 Crouzet, Nicolas, and Janice Eberly. 2023. “Rents and Intangible Capital: A Q+ Framework.” Journal of Finance 78 (4): 1873–1916.\par

% Reference 28; original PDF 55.
\noindent\hangindent=1.5em\hangafter=1 Curcuru, Stephanie E., Tomas Dvorak, and Francis E. Warnock. 2008. “Cross Border Returns Differentials.” Quarterly Journal of Economics 123 (4): 1495–1530.\par

% Reference 29; original PDF 55.
\noindent\hangindent=1.5em\hangafter=1 Curcuru, Stephanie E., Charles P. Thomas, and Francis E. Warnock. 2013. “On Returns Differentials.” Journal of International Money and Finance 36: 1–25.\par

% Reference 30; original PDF 55.
\noindent\hangindent=1.5em\hangafter=1 De Loecker, Jan, and Jan Eeckhout. 2021. “Global Market Power.” Unpublished.\par

% Reference 31; original PDF 55.
\noindent\hangindent=1.5em\hangafter=1 De Loecker, Jan, Jan Eeckhout, and Gabriel Unger. 2020. “The Rise of Market Power and the Macroeconomic Implications.” Quarterly Journal of Economics 135 (2): 561–644.\par

% Reference 32; original PDF 55.
\noindent\hangindent=1.5em\hangafter=1 Eggertsson, Gauti B., Jakob A. Robbins, and Ella Getz Wold. 2021. “Kaldor and Piketty’s Facts: The Rise of Monopoly Power in the United States.” Journal of Monetary Economics 124 (S): S19–S38.\par

% Reference 33; original PDF 55.
\noindent\hangindent=1.5em\hangafter=1 Engel, Charles, and John H. Rogers. 2006. “The US Current Account Deficit and the Expected Share of World Output.” Journal of Monetary Economics 53 (5): 1063–93.\par

% Reference 34; original PDF 55.
\noindent\hangindent=1.5em\hangafter=1 Farhi, Emmanuel, and François Gourio. 2018. “Accounting for Macro-Finance Trends: Market Power, Intangibles, and Risk Premia.” Brookings Papers on Economic Activity 49 (2): 147–250.\par

% Reference 35; original PDF 55.
\noindent\hangindent=1.5em\hangafter=1 Federal Reserve Bank of Philadelphia. 2024. Survey of Professional Forecasters, 10 Year Real GDP Growth Rate (RGDP10), Median Response. \url{https://www.philadelphiafed.org/surveys-and-data/rgdp10} (accessed March 2024).\par

% Reference 36; original PDF 55.
\noindent\hangindent=1.5em\hangafter=1 Gourinchas, Pierre-Olivier. 2023. “International Macroeconomics: From the Great Financial Crisis to COVID-19, and Beyond.” IMF Economic Review 71 (1): 1–34.\par

% Reference 37; original PDF 55.
\noindent\hangindent=1.5em\hangafter=1 Gourinchas, Pierre-Olivier, and Hélène Rey. 2007. “International Financial Adjustment.” Journal of Political Economy 115 (4): 665–703.\par

% Reference 38; original PDF 55.
\noindent\hangindent=1.5em\hangafter=1 Gourinchas, Pierre-Olivier, and Hélène Rey. 2014. “External Adjustment, Global Imbalances, Valuation Effects.” In Handbook of International Economics, Vol. 4, edited by Gita Gopinath, Elhanan Helpman, and Kenneth Rogoff, 585–645. North Holland.\par

% Reference 39; original PDF 55.
\noindent\hangindent=1.5em\hangafter=1 Gourinchas, Pierre-Olivier, Helene Rey, and Nicolas Govillot. 2017. “Exorbitant Privilege and Exorbitant Duty.” Unpublished.\par

% Reference 40; original PDF 55.
\noindent\hangindent=1.5em\hangafter=1 Greenwald, Daniel L., Martin Lettau, and Sydney Ludvigson. 2025. “How the Wealth Was Won: Factor Shares as Market Fundamentals.” Journal of Political Economy 133 (4): 1083–1132.\par

% Reference 41; original PDF 55.
\noindent\hangindent=1.5em\hangafter=1 Gutiérrez, Germán, and Thomas Philippon. 2017. “Investmentless Growth: An Empirical Investigation.” Brookings Papers on Economic Activity 48 (2): 89–169.\par



% ===== ORIGINAL PDF PAGE 56; JOURNAL PAGE 2206 =====
\sourcepage{56}{2206}

% Source PDF 56, text block 1.
% Reference 42; original PDF 56.
\noindent\hangindent=1.5em\hangafter=1 Gutierrez, German, and Thomas Philippon. 2023. “How European Markets Became Free: A Study of Institutional Drift.” Journal of the European Economic Association 21 (1): 251–92.\par

% Reference 43; original PDF 56.
\noindent\hangindent=1.5em\hangafter=1 Gutiérrez, Germán, and Sophie Piton. 2020. “Revisiting the Global Decline of the (Non-housing) Labor Share.” American Economic Review: Insights 2 (3): 321–38.\par

% Reference 44; original PDF 56.
\noindent\hangindent=1.5em\hangafter=1 Guvenen, Fatih, Raymond J. Mataloni Jr., Dylan G. Rassier, and Kim J. Ruhl. 2022. “Offshore Profit Shifting and Aggregate Measurement: Balance of Payments, Foreign Investment, Productivity, and the Labor Share.” American Economic Review 112 (6): 1848–84.\par

% Reference 45; original PDF 56.
\noindent\hangindent=1.5em\hangafter=1 Hall, Robert E. 2001. “The Stock Market and Capital Accumulation.” American Economic Review 91 (5): 1185–1202.\par

% Reference 46; original PDF 56.
\noindent\hangindent=1.5em\hangafter=1 Heathcote, Jonathan, and Fabrizio Perri. 2013. “The International Diversification Puzzle Is Not as Bad as You Think.” Journal of Political Economy 121 (6): 1108–59.\par

% Reference 47; original PDF 56.
\noindent\hangindent=1.5em\hangafter=1 Hsieh, Chang-Tai, and Peter J. Klenow. 2009. “Misallocation and Manufacturing TFP in China and India.” Quarterly Journal of Economics 124 (4): 1403–48.\par

% Reference 48; original PDF 56.
\noindent\hangindent=1.5em\hangafter=1 Jiang, Zhengyang, Robert J. Richmond, and Tony Zhang. 2024. “A Portfolio Approach to Global Imbalances.” Journal of Finance 79 (3): 2025–76.\par

% Reference 49; original PDF 56.
\noindent\hangindent=1.5em\hangafter=1 Karabarbounis, Loukas, and Brent Neiman. 2019. “Accounting for Factorless Income.” NBER Macroeconomics Annual 33: 167–228.\par

% Reference 50; original PDF 56.
\noindent\hangindent=1.5em\hangafter=1 Lane, Phillip R. 2020. “The Analytical Contribution of External Statistics: Addressing the Challenges.” Keynote Speech, Joint European Central Bank, Irving Fisher Committee and Banco de Portugal Conference on Bridging Measurement Challenges and Analytical Needs of External Statistics, Lisbon, February 17.\par

% Reference 51; original PDF 56.
\noindent\hangindent=1.5em\hangafter=1 Lane, Phillip R., and Gian Maria Milesi-Ferretti. 2001. “The External Wealth of Nations: Measures of Foreign Assets and Liabilities for Industrial and Developing Economies.” Journal of International Economics 55 (2): 263–94.\par

% Reference 52; original PDF 56.
\noindent\hangindent=1.5em\hangafter=1 Lane, Phillip R., and Gian Maria Milesi-Ferretti. 2007. “The External Wealth of Nations Mark II: Revised and Extended Estimates of Foreign Assets and Liabilities, 1970–2004.” Journal of International Economics 73 (2): 223–50.\par

% Reference 53; original PDF 56.
\noindent\hangindent=1.5em\hangafter=1 Lequiller, François, and Derek Blades. 2014. Understanding National Accounts, Second Edition. OECD Publishing.\par

% Reference 54; original PDF 56.
\noindent\hangindent=1.5em\hangafter=1 Lustig, Hanno, and Stijn Van Nieuwerburgh. 2008. “The Returns on Human Capital: Good News on Wall Street Is Bad News on Main Street.” Review of Financial Studies 21 (5): 2097–2137.\par

% Reference 55; original PDF 56.
\noindent\hangindent=1.5em\hangafter=1 Maggiori, Matteo, Brent Neiman, and Jesse Schreger. 2020. “International Currencies and Capital Allocation.” Journal of Political Economy 128 (6): 2019–66.\par

% Reference 56; original PDF 56.
\noindent\hangindent=1.5em\hangafter=1 Milesi-Ferretti, Gian Maria. 2021. “The US Is Increasingly a Net Debtor Nation. Should We Worry?” Brookings Institution, April 14. \url{https://www.brookings.edu/articles/the-us-is-increasingly-a-net-debtor-nation-should-we-worry/}.\par

% Reference 57; original PDF 56.
\noindent\hangindent=1.5em\hangafter=1 Milesi-Ferretti, Gian Maria. 2023. “Many Creditors, One Large Debtor: Understanding the Buildup of Global Stock Imbalances after the Global Financial Crisis.” Hutchins Center at Brookings Working Paper 90.\par

% Reference 58; original PDF 56.
\noindent\hangindent=1.5em\hangafter=1 MSCI, Inc. 2024. “MSCI Data, Developed Markets Standard (Price).” \url{https://www.msci.com/end-of-day-data-search} (accessed March 2024).\par

% Reference 59; original PDF 56.
\noindent\hangindent=1.5em\hangafter=1 Obstfeld, Maurice, and Kenneth Rogoff. 1995. “The Intertemporal Approach to the Current Account.” In Handbook of International Economics, Vol. 3, edited by Gene M. Grossman and Kenneth Rogoff, 1731–99. Elsevier.\par

% Reference 60; original PDF 56.
\noindent\hangindent=1.5em\hangafter=1 OECD. 2024. OECD Data Explorer. \url{https://data-explorer.oecd.org/} (accessed March 2024).\par

% Reference 61; original PDF 56.
\noindent\hangindent=1.5em\hangafter=1 Peters, Michael. 2020. “Heterogeneous Markups, Growth, and Endogenous Misallocation.” Econometrica 88 (5): 2037–73.\par

% Reference 62; original PDF 56.
\noindent\hangindent=1.5em\hangafter=1 Philippon, Thomas. 2019. The Great Reversal: How America Gave Up on Free Markets. Harvard University Press.\par

% Reference 63; original PDF 56.
\noindent\hangindent=1.5em\hangafter=1 Zucman, Gabriel. 2013. “The Missing Wealth of Nations: Are Europe and the US Net Debtors or Net Creditors?” Quarterly Journal of Economics 128 (3): 1321–64.\par


\end{document}
